% !TEX program = xelatex
% ============================================================
% مقاله‌ی پژوهشی پزشکی - بستری مجدد ظرف ۳۰ روز در بیماران دیابتی
% (ترجمه‌ی فارسی؛ کامپایل با XeLaTeX)
% ------------------------------------------------------------
% توضیح فنی: این فایل به بسته‌های bidi / polyglossia / xepersian
% نیازی ندارد. راست‌به‌چپ‌نویسی مستقیماً با دستورهای TeX--XeT
% (\beginR ... \endR و \beginL ... \endL) پیاده‌سازی شده است.
% فونت‌ها: بی‌زر (متن)، بی‌تیتر (تیترها)، DejaVu Serif / DejaVu Sans Mono (لاتین).
% اعداد لاتین در فونت فارسی به‌طور خودکار به ارقام فارسی نگاشت می‌شوند.
% دستور \lr{...} برای متن/اعداد چپ‌به‌راست (لاتین) است.
% ============================================================
\documentclass[11pt,a4paper]{article}

\usepackage[a4paper,margin=2.2cm,top=2.3cm,bottom=2.4cm]{geometry}
\usepackage[no-math]{fontspec}
\TeXXeTstate=1

% ---------- fonts ----------
% متن: «بی‌زر» (B Zar)؛ تیترها: «بی‌تیتر» (B Titr). فایل‌های فونت در پوشه‌ی fonts/ قرار دارند.
\setmainfont[Path=fonts/,Extension=.ttf,UprightFont=BZar,BoldFont=BZarBold,
  ItalicFont=BZar,BoldItalicFont=BZarBold,
  Script=Arabic,Mapping=farsidigits,Scale=1.0]{BZar}
\newfontfamily\titrfont[Path=fonts/,Extension=.ttf,UprightFont=BTitrBold,BoldFont=BTitrBold,
  Script=Arabic,Mapping=farsidigits,Scale=1.0]{BTitrBold}
\newfontfamily\latinfont[Scale=0.9,Ligatures=TeX]{DejaVu Serif}
\setsansfont[Scale=0.9]{DejaVu Sans}
\setmonofont[Scale=0.82]{DejaVu Sans Mono}
\DeclareRobustCommand{\lr}[1]{\beginL{\latinfont ^^^^200e#1^^^^200e}\endL}
\DeclareRobustCommand{\code}[1]{\beginL{\texttt{#1}}\endL}
\newcommand{\usc}{\mathord{\rule[-0.35ex]{0.5em}{0.4pt}}}
\newcommand{\lessthirty}{\ensuremath{\mathtt{<30}}}
\newcommand{\gtthirty}{\ensuremath{\mathtt{>30}}}
\newcommand{\po}{)}% پرانتز بازکننده‌ی چسبیده به \lr (شکل آینه‌ای)
\newcommand{\pc}{(}% پرانتز بسته‌کننده‌ی چسبیده به \lr (شکل آینه‌ای)
\newcommand{\fat}[1]{\text{\beginR #1\endR}}
\DeclareRobustCommand{\emph}[1]{\textbf{#1}}% فارسی ایتالیک ندارد

\usepackage{setspace}
\setstretch{1.32}
\setlength{\emergencystretch}{4em}
\sloppy

\usepackage{amsmath,amssymb}
\usepackage{booktabs,longtable,tabularx,array}
\usepackage{graphicx}
\usepackage{float}
\usepackage{capt-of}
\usepackage{xcolor}
\definecolor{medicalblue}{HTML}{174A6E}
\definecolor{medicallight}{HTML}{EAF2F7}
\definecolor{darkgray}{HTML}{3D4650}
\usepackage{fancyhdr}
\usepackage{lastpage}
\usepackage{url}
\usepackage[hidelinks,unicode,pdftitle={پیش‌بینی بستری مجدد ظرف ۳۰ روز در بیماران مبتلا به دیابت},pdfauthor={Atefe Asadi, Iman Jamshidi, Mahdiyeh Mirzaei Valojerdie}]{hyperref}

% ---------- names ----------
\renewcommand{\figurename}{شکل}
\renewcommand{\tablename}{جدول}
\renewcommand{\contentsname}{فهرست مطالب}
\renewcommand{\listfigurename}{فهرست شکل‌ها}
\renewcommand{\listtablename}{فهرست جدول‌ها}
\renewcommand{\abstractname}{چکیده}
\renewcommand{\appendixname}{پیوست}

\makeatletter
% ---------- RTL paragraph machinery ----------
\newskip\fa@indent
\fa@indent=1.4em\relax
\let\faindent\fa@indent
\def\fa@ep{\beginR\hskip\fa@indent\relax}
\parindent=0pt
\parskip=0.15\baselineskip plus 1pt
\AddToHook{para/begin}{\fa@ep}
\def\@hangfrom#1{\setbox\@tempboxa\hbox{{#1}}\hangindent-\wd\@tempboxa\noindent\box\@tempboxa}

% ---------- sectioning ----------
\renewcommand\thesection{\arabic{section}}
\renewcommand\@seccntformat[1]{\csname the#1\endcsname\hskip.7em}
\renewcommand\section{\@startsection{section}{1}{\z@}%
  {-3.2ex \@plus -1ex \@minus -.2ex}{1.8ex \@plus.2ex}%
  {\titrfont\Large\color{medicalblue}\fa@indent=0pt\relax}}
\renewcommand\subsection{\@startsection{subsection}{2}{\z@}%
  {-2.6ex\@plus -1ex \@minus -.2ex}{1.2ex \@plus .2ex}%
  {\titrfont\large\color{darkgray}\fa@indent=0pt\relax}}
\renewcommand\subsubsection{\@startsection{subsubsection}{3}{\z@}%
  {-2.2ex\@plus -1ex \@minus -.2ex}{1ex \@plus .2ex}%
  {\titrfont\normalsize\color{darkgray}\fa@indent=0pt\relax}}
\newcommand{\faalph}[1]{\ifcase#1\or الف\or ب\or پ\or ت\fi}
\renewcommand\appendix{\par\setcounter{section}{0}\setcounter{subsection}{0}%
  \gdef\thesection{\faalph{\value{section}}}}

% ---------- lists (custom RTL) ----------
\newcounter{faitem}
\newdimen\fa@lw
\fa@lw=1.2cm
\newenvironment{faenum}
  {\par\vskip.25\baselineskip\setcounter{faitem}{0}\fa@indent=0pt\relax
   \renewcommand\faitem{\par\refstepcounter{faitem}\hangindent=-\fa@lw\hangafter=1
     \noindent\hbox to\fa@lw{\beginR\thefaitem.\hfil\endR}}}
  {\par\vskip.25\baselineskip}
\newenvironment{faitemize}
  {\par\vskip.25\baselineskip\fa@indent=0pt\relax
   \renewcommand\faitem{\par\hangindent=-\fa@lw\hangafter=1
     \noindent\hbox to\fa@lw{\beginR\ensuremath{\bullet}\hfil\endR}}}
  {\par\vskip.25\baselineskip}
\newcommand\faitem{}

% ---------- display math (kept LTR, centred) ----------
\newcommand{\mathblock}[1]{\par\vskip.5\baselineskip
  \hbox to\hsize{\hss$\displaystyle #1$\hss}\vskip.5\baselineskip}

% ---------- captions ----------
\setlength\abovecaptionskip{8\p@}
\setlength\belowcaptionskip{4\p@}
\long\def\@makecaption#1#2{%
  \vskip\abovecaptionskip
  {\small\fa@indent=0pt\relax
   \sbox\@tempboxa{\beginR\textbf{#1:}\ #2\endR}%
   \ifdim\wd\@tempboxa>\hsize
     \noindent\textbf{#1:}\ #2\par
   \else
     \global\@minipagefalse
     \hb@xt@\hsize{\hfil\box\@tempboxa\hfil}%
   \fi}%
  \vskip\belowcaptionskip}

\def\LT@makecaption#1#2#3{\LT@mcol\LT@cols c{\hbox to\z@{\hss\parbox[t]\LTcapwidth{\small\faindent=0pt\relax\centering\textbf{#2:}\ #3\endgraf\vskip\baselineskip}\hss}}}

% ---------- figures (placeholder when image file is missing) ----------
\newcommand{\fig}[2]{\begin{figure}[H]\centering
  \IfFileExists{Figures/#2}%
    {\hbox to\hsize{\hfil\includegraphics[width=#1\textwidth]{Figures/#2}\hfil}}%
    {\hbox to\hsize{\hfil\fbox{\parbox[c][3.2cm][c]{#1\textwidth}{\centering\fa@indent=0pt\relax\small
      [محل شکل - فایل تصویر «\lr{#2}» در پوشه‌ی \lr{Figures} یافت نشد]}}\hfil}}%
  \end{figure}}

% ---------- table of contents (RTL) ----------
\renewcommand*\numberline[1]{\hb@xt@\@tempdima{\hfil #1\hskip .6em}}
\renewcommand*\@dottedtocline[5]{%
  \ifnum #1>\c@tocdepth \else
    \vskip \z@ \@plus.2\p@
    {\fa@indent=#2\relax
     \leftskip \@tocrmarg \parfillskip -\leftskip
     \interlinepenalty\@M
     \@tempdima #3\relax
     \hangindent=-\dimexpr#2+#3\relax \hangafter=1
     \noindent
     {#4}\nobreak
     \leaders\hbox{$\m@th \mkern \@dotsep mu\hbox{.}\mkern \@dotsep mu$}\hfill
     \nobreak
     \hb@xt@\@pnumwidth{\normalfont \normalcolor #5\hfil}%
     \par}%
  \fi}

% ---------- header / footer ----------
\pagestyle{fancy}
\fancyhf{}
\fancyhead[R]{\small\hbox to 10cm{\beginR بستری مجدد ظرف ۳۰ روز در بیماران مبتلا به دیابت\hfil\endR}}
\fancyhead[L]{\small\hbox to 6cm{\beginR\hfil گزارش علم داده‌های پزشکی\endR}}
\fancyfoot[C]{\small\faindent=0pt\beginR صفحه \thepage\ از \pageref{LastPage}\endR}
\renewcommand{\headrulewidth}{0.4pt}
\setlength{\headheight}{14pt}

% ---------- abstract ----------
\newenvironment{faabstract}
  {\par\begingroup\small\leftskip=1cm\rightskip=1cm
   {\centering\titrfont\large\color{medicalblue}\fa@indent=0pt\relax\abstractname\par}\vskip.6em}
  {\par\endgroup}

\setcounter{tocdepth}{2}
\setcounter{secnumdepth}{3}
\hypersetup{pdfencoding=unicode}
\raggedbottom
\makeatother

\begin{document}

% ============================================================
% صفحه‌ی عنوان
% ============================================================
\begin{titlepage}
\centering\faindent=0pt\relax
\vspace*{0.6cm}
{\color{medicalblue}\titrfont\Large تحلیل داده‌محور پزشکی\par}\vspace{0.6cm}
{\titrfont\Huge پیش‌بینی بستری مجدد ظرف ۳۰ روز در بیماران مبتلا به دیابت\par}\vspace{0.5cm}
{\Large تحلیل آماری و یادگیری ماشین آگاه از نشت داده بر روی ۱۰۱,۷۶۶ مراجعه‌ی بیمارستانی\par}\vspace{0.8cm}
{\large\bfseries عاطفه اسدی\par}\vspace{0.05cm}
{\large\bfseries ایمان جمشیدی\par}\vspace{0.05cm}
{\large\bfseries مهدیه میرزایی ولوجردی\par}\vspace{0.15cm}
{\small\lr{Atefe Asadi \quad$\cdot$\quad Iman Jamshidi \quad$\cdot$\quad Mahdiyeh Mirzaei Valojerdie}\par}\vspace{0.8cm}

\setlength{\fboxsep}{9pt}\setlength{\fboxrule}{0.8pt}
\hbox to\hsize{\hfil\fcolorbox{medicalblue}{medicallight}{\parbox{0.84\textwidth}{%
\centering\faindent=0pt\relax
\textbf{مقاله‌ی پژوهشی / گزارش علم داده‌های پزشکی}\\[0.25cm]
تحلیل هم‌گروهی مشاهده‌ای \quad$\cdot$\quad آماره‌های اندازه‌ی اثر \quad$\cdot$\quad یادگیری ماشین\\
اعتبارسنجی گروه‌بندی‌شده بر اساس بیمار \quad$\cdot$\quad تحلیل خطا در زیرگروه‌ها}}\hfil}

\vfill
\small
\begin{tabular}{@{}>{\raggedleft\arraybackslash}p{9.4cm}>{\beginR}r<{\endR}@{}}
\beginR مجموعه‌داده‌ی «دیابت، ۱۳۰ بیمارستان آمریکا، سال‌های ۱۹۹۹ تا ۲۰۰۸»\endR & \textbf{مجموعه‌داده:}\\[2pt]
\beginR بستری مجدد ظرف ۳۰ روز \po\lessthirty\pc\endR & \textbf{پیامد اصلی:}\\[2pt]
\beginR ۱۰۱,۷۶۶ مراجعه‌ی بیمارستانی\endR & \textbf{جامعه‌ی مورد مطالعه:}\\[2pt]
\beginR ۹۹,۳۴۳ مراجعه پس از حذف پرونده‌های فوت/هاسپیس\endR & \textbf{هم‌گروه مدل‌سازی:}\\[2pt]
\beginR رگرسیون لجستیک، درخت تصمیم، جنگل تصادفی، \lr{XGBoost}\endR & \textbf{چارچوب اصلی مدل‌سازی:}\\[2pt]
\beginR \lr{Python, pandas, SciPy, scikit-learn, XGBoost, Matplotlib, Seaborn}\endR & \textbf{نرم‌افزار:}\\
\end{tabular}

\vfill
{\large\textbf{تهیه‌شده به‌صورت دست‌نوشته‌ی تحلیلی در قالب مقاله‌ی علمی}}\\
\vspace{0.3cm}
{\small نسخه‌ی تهیه‌شده از آرشیو پروژه‌ی ارائه‌شده}\\
{\small اکتبر ۲۰۲۶ (مهر ۱۴۰۵)}
\end{titlepage}

% ============================================================
% چکیده
% ============================================================
\newcommand{\fapagealpha}[1]{\ifcase#1\or الف\or ب\or پ\or ت\or ث\or ج\or چ\or ح\or خ\or د\or ذ\or ر\else #1\fi}
\pagenumbering{arabic}\renewcommand{\thepage}{\fapagealpha{\value{page}}}
\begin{faabstract}
\textbf{زمینه:} بستری مجدد ظرف ۳۰ روز، پیامدی مهم برای سنجش کیفیت مراقبت و هماهنگی خدمات درمانی است، به‌ویژه در بیماران مبتلا به بیماری‌های مزمن پیچیده. مراجعه‌های بیمارستانی مرتبط با دیابت با تشخیص‌های ناهمگن، بیماری‌های همراه، مواجهه‌ی دارویی، سابقه‌ی استفاده از خدمات درمانی و مسیرهای متنوع ترخیص همراه‌اند. مدل‌های پیش‌بینی می‌توانند به طبقه‌بندی ریسک کمک کنند، اما عدم توازن کلاس‌ها، مراجعه‌های مکرر یک بیمار، مصنوعات کدگذاری و ناپایداری عملکرد در زیرگروه‌ها ممکن است نتایجی گمراه‌کننده و بیش‌ازحد خوش‌بینانه به بار آورند.

\textbf{هدف:} توصیف عوامل مرتبط با بستری مجدد زودهنگام و ارزیابی مدل‌های یادگیری ماشین در سطح بیمار و آگاه از نشت داده، برای شناسایی مراجعه‌هایی که ظرف ۳۰ روز با بستری مجدد دنبال شده‌اند.

\textbf{روش‌ها:} این تحلیل از ۱۰۱,۷۶۶ مراجعه‌ی موجود در مجموعه‌داده‌ی «دیابت، ۱۳۰ بیمارستان آمریکا» استفاده کرد. تحلیل اکتشافی شامل شمارش مراجعه‌ها، نرخ‌های ویژه‌ی هر پیامد، تفاوت‌های میانگین استانداردشده \lr{(SMD)}، ضریب \lr{V} کرامر، همبستگی پیرسون و غربالگری با دامنه‌ی بین‌چارکی بود. برای پیش‌بینی، پس از حذف ترخیص‌های منجر به فوت یا هاسپیس، ۹۹,۳۴۳ مراجعه باقی ماند. با \lr{GroupShuffleSplit} بر اساس شناسه‌ی بیمار، یک تقسیم آموزش/آزمون ۸۰/۲۰ بدون هیچ هم‌پوشانی بیماری ساخته شد. مدل‌های رگرسیون لجستیک، درخت تصمیم، جنگل تصادفی و \lr{XGBoost} در شرایط عدم توازن کلاس‌ها ارزیابی شدند. آستانه‌های بازیابی‌محور با امتیاز \lr{$F_2$} در مجموعه‌ی اعتبارسنجی انتخاب شد. یک مدل \lr{XGBoost} مرحله‌ی دوم با اعتبارسنجی متقاطع تصادفی گروه‌بندی‌شده بهینه شد. عملکرد در زیرگروه‌ها با بازیابی \lr{(Recall)} و نرخ منفی کاذب \lr{(FNR)} در نژاد، جنسیت و گروه‌های سنی بررسی شد.

\textbf{یافته‌ها:} بستری مجدد زودهنگام در ۱۱/۱۶\% از کل مراجعه‌ها (۱۱,۳۵۷ از ۱۰۱,۷۶۶) و در ۱۱/۳۹\% از هم‌گروه مدل‌سازی رخ داد. سابقه‌ی بستری قبلی روشن‌ترین سیگنال عددی بود: میانگین بستری‌های قبلی در بستری‌های مجدد زودهنگام ۱/۲۲ و در مراجعه‌های بدون بستری مجدد زودهنگام ۰/۳۸ بود \lr{(SMD = 0.415)}. مقادیر \lr{V} کرامر در مجموع کوچک بودند و مقصد ترخیص بیشترین ارتباط طبقه‌ای را داشت \lr{(V = 0.125)}. مدل \lr{XGBoost} مرحله‌ی الف به \lr{ROC--AUC = 0.683} و \lr{PR--AUC = 0.238} رسید. جایگزینی آستانه‌ی پیش‌فرض ۰/۵۰ با آستانه‌ی ۰/۴۱ که روی مجموعه‌ی اعتبارسنجی تنظیم شد، بازیابی \lr{XGBoost} را از ۰/۵۹۸ به ۰/۷۹۷ رساند، در حالی که \lr{ROC--AUC} و \lr{PR--AUC} تغییری نکردند؛ بهای این افزایش، کاهش دقت \lr{(Precision)} از ۰/۱۸۸ به ۰/۱۵۲ و کاهش صحت \lr{(Accuracy)} از ۰/۶۶۳ به ۰/۴۷۴ بود و حدود ۵۹\% از مراجعه‌های مجموعه‌ی آزمون علامت‌گذاری شدند. مدل نهایی \lr{XGBoost} با ابرپارامترهای جست‌وجوشده، در آستانه‌ی ۰/۵۰ به بازیابی ۰/۶۱۹، دقت ۰/۱۸۱، \lr{F1} برابر ۰/۲۸۰ و \lr{ROC--AUC} نزدیک به ۰/۶۸ دست یافت. شکاف \lr{FNR} میان گروه‌های سنی در هر چهار مدل شدید بود (از ۰/۸۸۵ تا ۰/۹۶۷، و \lr{FNR = 1.00} دست‌کم در یک گروه سنی)؛ جوان‌ترین گروه‌های سنی درگیر فقط ۱ و ۶ بستری مجدد زودهنگام داشتند، بنابراین بزرگی دقیق این شکاف نامطمئن است، اما نمی‌توان مدل‌ها را برای این بیماران بدون اعتبارسنجی جداگانه ایمن دانست.

\textbf{نتیجه‌گیری:} سابقه‌ی بستری قبلی و مقصد ترخیص، پایدارترین سیگنال‌های پیش‌بینی‌کننده بودند، در حالی که توان تفکیک کلی در حد متوسط ماند. یافته‌ها از به‌کارگیری مدل به‌عنوان ابزار کمکی طبقه‌بندی ریسک با بازبینی انسانی پشتیبانی می‌کنند، نه تصمیم‌گیری بالینی خودکار. پیش از هرگونه استقرار بالینی، اعتبارسنجی بیرونی و به‌روز، ارزیابی کالیبراسیون، بازه‌های عدم قطعیت و نمونه‌های بزرگ‌تر در زیرگروه‌ها ضروری است.
\end{faabstract}

\medskip
\noindent\textbf{واژه‌های کلیدی:} دیابت؛ بستری مجدد بیمارستانی؛ بستری مجدد ظرف ۳۰ روز؛ تحلیل داده‌های سلامت؛ یادگیری ماشین؛ \lr{XGBoost}؛ طبقه‌بندی نامتوازن؛ آستانه‌ی تصمیم؛ نشت داده در سطح بیمار؛ انصاف الگوریتمی؛ طبقه‌بندی ریسک.

\clearpage
\tableofcontents
\clearpage
\listoffigures
\clearpage
\listoftables
\clearpage
\pagenumbering{arabic}\renewcommand{\thepage}{\arabic{page}}

% ============================================================
\section{مقدمه}
% ============================================================

بستری مجدد بیمارستانی پیامدی مهم از نظر بالینی و عملیاتی است، زیرا ترکیبی از بار بیماری، سابقه‌ی استفاده از خدمات درمانی، شرایط ترخیص، تداوم مراقبت و سیر طبیعی بیماری مزمن را بازتاب می‌دهد. در بیماران مبتلا به دیابت، این عوامل می‌توانند با بیماری‌های همراه متعدد، تغییرات دارویی، عوارض حاد و مواجهه‌ی مکرر با خدمات بستری و اورژانس برهم‌کنش داشته باشند. در نتیجه، یک مراجعه‌ی اداری بیمارستانی ممکن است تنها یک مشاهده در الگویی طولی از مراقبت باشد.

پژوهش حاضر مسئله‌ی شناسایی مراجعه‌های بیمارستانی را بررسی می‌کند که ظرف ۳۰ روز با بستری مجدد دنبال شده‌اند. تحلیل بر پایه‌ی مجموعه‌داده‌ی «دیابت، ۱۳۰ بیمارستان آمریکا» انجام شده است که حدود صد هزار مراجعه‌ی بیمارستانی بیماران مبتلا به دیابت را در چندین بیمارستان آمریکا طی سال‌های ۱۹۹۹ تا ۲۰۰۸ در بر می‌گیرد. این مجموعه‌داده برای کارهای روش‌شناختی به‌ویژه سودمند است، زیرا متغیرهای بالینی و اداری را با هم دارد، شامل مراجعه‌های مکرر یک بیمار است، مقادیر گمشده‌ی قابل‌توجهی دارد و پیامد بستری مجدد زودهنگام در آن به‌شدت نامتوازن است.

چالش روش‌شناختی اصلی این است که اگر مراجعه‌های مکرر یک بیمار به‌صورت تصادفی میان مجموعه‌ی آموزش و آزمون تقسیم شوند، نشت اطلاعات رخ می‌دهد. در این حالت، مدل ممکن است خوب تعمیم‌پذیر به نظر برسد، در حالی که بخشی از آن صرفاً الگوهای اختصاصی بیمار را می‌شناسد که در هر دو بخش حضور دارند. بنابراین این پژوهش از تقسیم‌بندی گروهی بر اساس بیمار و اعتبارسنجی متقاطع گروهی استفاده می‌کند. هدف، صرفاً بیشینه‌کردن یک معیار متعارف صحت نیست؛ بلکه بررسی این است که آیا مدل‌ها می‌توانند نسبت معناداری از موارد بستری مجدد زودهنگام را شناسایی کنند، در حالی که مصالحه‌ی میان بازیابی و دقت به‌روشنی نشان داده می‌شود.

این کار چهار هدف مکمل دارد:
\begin{faenum}
\faitem توصیف توزیع بستری مجدد ظرف ۳۰ روز و الگوهای مهم استفاده‌ی قبلی از خدمات؛
\faitem کمّی‌سازی تفاوت‌ها با اندازه‌های اثرِ تفسیرپذیر، به‌جای تکیه‌ی صرف بر معناداری آماری؛
\faitem مقایسه‌ی چهار الگوریتم طبقه‌بندی تحت کنترل نشت در سطح بیمار و عدم توازن کلاس‌ها؛
\faitem بررسی الگوهای خطا در زیرگروه‌ها و برگرداندن یافته‌ها به پیامدهای محتاطانه و بالینی‌مرتبط برای پشتیبانی از تصمیم.
\end{faenum}

تمایز میان ارتباط و علیت در سراسر متن حفظ شده است. ویژگی‌ای که ارتباط قوی با بستری مجدد دارد لزوماً علتِ قابل‌مداخله‌ی بستری مجدد نیست، و اهمیت یک ویژگی در مدل، شاهدی بر علیت بالینی نیست.

% ============================================================
\section{زمینه‌ی بالینی و تحلیلی}
% ============================================================

\subsection{بستری مجدد ظرف ۳۰ روز به‌عنوان یک پیامد پزشکی}

بستری مجدد ظرف ۳۰ روز معمولاً به‌عنوان پیامدی کوتاه‌افق برای ارزیابی فرایندهای مراقبت پس از ترخیص و شناسایی بیمارانی به کار می‌رود که ممکن است به پیگیری نزدیک‌تر نیاز داشته باشند. در بیماری‌های مزمن مانند دیابت، بستری مجدد می‌تواند ناشی از عوارض عودکننده‌ی بیماری، بیماری‌های همراه، برنامه‌ریزی ناکافی برای انتقال مراقبت، مشکلات دارویی یا رویدادهای حاد نامرتبط باشد. از آنجا که مجموعه‌داده‌ی حاضر مشاهده‌ای و اداری است، این تحلیل تلاشی برای نسبت‌دادن هر بستری مجدد به سازوکار بالینی مشخص نمی‌کند.

هدف (متغیر پاسخ) به این صورت تعریف می‌شود:
\mathblock{Y=\begin{cases}1, & \fat{اگر مراجعه ظرف ۳۰ روز با بستری مجدد دنبال شود،}\\[3pt] 0, & \fat{در غیر این صورت.}\end{cases}}
بنابراین پیامد سه‌سطحی اصلی \po{}\lessthirty{}، \gtthirty{} و \code{NO}\pc{} برای مدل‌سازی به یک پیامد دودویی تبدیل می‌شود. این صورت‌بندی مستقیماً پرسش بالینی «عود زودهنگام پس از ترخیص» را بازتاب می‌دهد.

\subsection{چرا عدم توازن کلاس‌ها اهمیت دارد}

تنها ۱۱/۱۶\% از کل مجموعه‌داده به کلاس بستری مجدد زودهنگام تعلق دارد. طبقه‌بندی‌ای که همه‌ی مراجعه‌ها را منفی برچسب بزند، صحتی ظاهری در حدود ۸۸/۸۴\% خواهد داشت، در حالی که هیچ‌یک از بیمارانی را که دچار بستری مجدد زودهنگام می‌شوند شناسایی نکرده است. از این رو صحت \lr{(Accuracy)} در این زمینه معیار سرخط مناسبی نیست.

برای یک طبقه‌بند دودویی، کمیت‌های اصلی عبارت‌اند از بازیابی \lr{(Recall)} و دقت \lr{(Precision)}:
\mathblock{\mathrm{Recall}=\frac{TP}{TP+FN},\qquad \mathrm{Precision}=\frac{TP}{TP+FP},}
و امتیاز \lr{$F_1$}:
\mathblock{F_1=\frac{2\,\mathrm{Precision}\,\mathrm{Recall}}{\mathrm{Precision}+\mathrm{Recall}}.}
هنگامی که از دست رفتن یک رویداد واقعی پرهزینه‌تر از تولید یک هشدار اضافی تلقی شود، می‌توان از معیار \lr{$F_\beta$} با $\beta>1$ استفاده کرد:
\mathblock{F_\beta=(1+\beta^2)\frac{\mathrm{Precision}\times\mathrm{Recall}}{\beta^2\,\mathrm{Precision}+\mathrm{Recall}}.}
این پژوهش $\beta=2$ را به کار می‌برد، که نتیجه می‌دهد:
\mathblock{F_2=5\,\frac{\mathrm{Precision}\times\mathrm{Recall}}{4\,\mathrm{Precision}+\mathrm{Recall}}.}

\subsection{توان تفکیک و رتبه‌بندی}

مساحت زیر منحنی \lr{ROC} \lr{(ROC--AUC)} توانایی مدل را در رتبه‌بندی مراجعه‌های مثبت بالاتر از مراجعه‌های منفی، در تمام آستانه‌های ممکن خلاصه می‌کند. تحلیل دقت-بازیابی \lr{(Precision--Recall)} وقتی کلاس مثبت نادر است، به‌ویژه آگاهی‌بخش است. بنابراین مساحت زیر منحنی دقت-بازیابی \lr{(PR--AUC)} مکمل \lr{ROC--AUC} است و کمک می‌کند معلوم شود مدل در شیوع پایین کلاس مثبت چقدر کارآمد می‌ماند.

نرخ منفی کاذب \lr{(FNR)} برابر است با:
\mathblock{\mathrm{FNR}=\frac{FN}{FN+TP}=1-\mathrm{Recall}.}
در کاربرد اولویت‌بندی (تریاژ) بالینی، \lr{FNR} از نظر عملیاتی اهمیت دارد، زیرا یک منفی کاذب به مراجعه‌ای از نوع بستری مجدد زودهنگام اشاره دارد که مدل آن را علامت‌گذاری نکرده است.

% ============================================================
\section{مواد و روش‌ها}
% ============================================================

\subsection{منبع داده و هم‌گروه مطالعه}

این پروژه از مجموعه‌داده‌ی «دیابت، ۱۳۰ بیمارستان آمریکا، سال‌های ۱۹۹۹ تا ۲۰۰۸» استفاده می‌کند. مجموعه‌داده‌ی اصلی شامل ۱۰۱,۷۶۶ مراجعه و ۷۱,۵۱۸ بیمار یکتا است و پس از پیش‌پردازش و بازبینی کیفیت داده در پروژه، ۵۰ متغیر دارد. داده‌ها نمایانگر مراجعه‌های بیمارستانی‌اند، نه افراد مستقل: ۵۴,۷۴۵ بیمار فقط یک مراجعه دارند، ۱۶,۷۷۳ بیمار بیش از یک مراجعه دارند و بیشینه‌ی تعداد مراجعه برای یک بیمار ۴۰ است. این تمایز در برآورد عملکرد مدل ضروری است.

پروژه یک مجموعه‌داده‌ی اکتشافیِ هم‌گروه کامل را برای تحلیل توصیفی نگه داشت. برای پیش‌بینی، ۲,۴۲۳ مراجعه‌ی منجر به فوت/هاسپیس حذف شد و ۹۹,۳۴۳ مراجعه باقی ماند. این حذف انجام شد زیرا شرایط ترخیص مربوط به آن‌ها مسئله‌ی متعارف ریسک بستری مجدد پس از ترخیص را نمایندگی نمی‌کند، و نیز به این دلیل که نشانگر \code{death\_or\_hospice} پس از این محدودسازی ثابت می‌شد.

\begin{table}[H]
\centering
\caption{جامعه‌ی مورد مطالعه و مراحل تحلیلی.}
\label{tab:cohort}
{\faindent=0pt\relax\begin{tabularx}{\textwidth}{>{\raggedleft\arraybackslash}X>{\beginR}r<{\endR}>{\beginR}r<{\endR}>{\beginR}r<{\endR}}
\toprule
هدف & بستری مجدد زودهنگام & مراجعه‌ها & مرحله\\

\midrule
تحلیل اکتشافی و آمار توصیفی & ۱۱,۳۵۷ (۱۱/۱۶\%) & ۱۰۱,۷۶۶ & کل مجموعه‌داده\\
مدل‌سازی پیش‌بینی & ۱۱/۳۹\% & ۹۹,۳۴۳ & هم‌گروه بدون فوت/هاسپیس\\
برازش مدل & -- & ۷۹,۳۹۵ & مجموعه‌ی آموزش\\
ارزیابی نهایی عملکرد & ۲,۲۶۰ (۱۱/۳۳\%) & ۱۹,۹۴۸ & مجموعه‌ی آزمون نگه‌داشته‌شده\\

\bottomrule
\end{tabularx}}
\end{table}

\subsection{کیفیت داده و پیش‌پردازش}

پروژه بررسی داده‌های تکراری، مقادیر گمشده، غربالگری متغیرهای ثابت، بازکدگذاری متغیرهای طبقه‌ای، استانداردسازی عددی و غربالگری داده‌های پرت را مستند کرده است. داده‌ها در متغیرهایی مانند \code{weight}، \code{payer\_code} و \code{medical\_specialty} مقادیر گمشده‌ی قابل‌توجهی دارند. فیلد \code{weight} حدود ۹۶/۸۶\% گمشده است و به‌عنوان اندازه‌گیری پیوسته‌ی قابل‌اعتماد در نظر گرفته نشد.

خط لوله‌ی پیش‌پردازش برای متغیرهای عددی از جای‌گذاری با میانه، و برای متغیرهای طبقه‌ای از پرتکرارترین مقدار یا یک دسته‌ی صریح «\lr{Missing}» استفاده می‌کند. متغیرهای عددی در مواردی که مدل لازم دارد استانداردسازی می‌شوند. متغیرهای طبقه‌ای درون یک خط لوله‌ی مدل‌سازی با روش یک‌داغ \lr{(one-hot)} کدگذاری می‌شوند تا تبدیل‌ها فقط از داده‌ی آموزش برآورد شوند.

دو محدودیت مهم پنهان نشد و نگه داشته شد. نخست، شناسه‌های اداری مانند \code{admission\_type\_id}، \code{discharge\_disposition\_id} و \code{admission\_source\_id} علاوه بر نسخه‌های کدگذاری‌شده‌ی طبقه‌ای‌شان، به‌طور خودکار عددی نیز در نظر گرفته شدند. دوم، ستون‌های خام تشخیص \code{diag\_1} تا \code{diag\_3} در کنار دسته‌های گروه‌بندی‌شده‌ی تشخیص باقی ماندند و فضای ویژگی یک‌داغِ بسیار پهنی ایجاد کردند. این انتخاب‌ها منابع بالقوه‌ی نوفه‌اند و به همین دلیل به‌عنوان محدودیت بحث می‌شوند و به‌عنوان سازوکار بالینی تفسیر نمی‌شوند.

\begin{table}[H]
\centering
\caption{ممیزی کیفیت داده و تصمیم‌های آماده‌سازی.}
\label{tab:dq}
\small
{\faindent=0pt\relax\begin{tabularx}{\textwidth}{>{\raggedleft\arraybackslash}X>{\beginR}r<{\endR}}
\toprule
نتیجه / تصمیم & گام\\

\midrule
۱۰۱,۷۶۶ سطر $\times$ ۵۰ ستون؛ فایل‌ها با \code{keep\_default\_na=False} خوانده شدند تا رشته‌ی لفظی \code{None} در \code{A1Cresult} و \code{max\_glu\_serum} به‌عنوان یک دسته‌ی معتبر حفظ شود؛ علامت \code{?} به‌عنوان نشانگر مقدار گمشده در نظر گرفته شد & بارگذاری\\
۰ شناسه‌ی مراجعه‌ی تکراری و ۰ سطر کاملاً تکراری؛ هیچ سطری حذف نشد & مراجعه‌های تکراری\\
۱۶,۷۷۳ بیمار با چند مراجعه حفظ شدند؛ شناسه‌ی بیمار برای تقسیم‌بندی گروهی به کار رفت & مراجعه‌های مکرر\\
\code{weight}: ۹۶/۸۶\% (کنار گذاشته شد)؛ \code{medical\_specialty}: ۴۹/۰۸\%؛ \code{payer\_code}: ۳۹/۵۶\%؛ \code{race}: ۲/۲۳\%؛ \code{diag\_3}: ۱/۴۰\%؛ \code{diag\_2}: ۰/۳۵\%؛ \code{diag\_1}: ۰/۰۲\%. موارد اخیر به یک دسته‌ی صریح \lr{Missing} بازکدگذاری شدند؛ هیچ سطری حذف نشد. سه مقدار جنسیتِ نامشخص/نامعتبر نیز به \lr{Missing} بازکدگذاری شد & مقادیر گمشده\\
\code{examide} و \code{citoglipton} فقط یک مقدار یکتا دارند و کنار گذاشته شدند & ویژگی‌های ثابت\\
حصارهای $1.5\times$\lr{IQR} مقادیر زیادی را علامت می‌زنند (۱۶/۴۵\% از \code{number\_outpatient} و ۱۱/۱۹\% از \code{number\_emergency})، اما این متغیرها صفر-متورم‌اند (به‌ترتیب ۸۳/۵۵\% و ۸۸/۸۱\% صفر)؛ \lr{IQR} فقط ابزار غربالگری بود و هیچ مشاهده‌ای حذف نشد & غربالگری داده‌های پرت\\
\code{admission\_type\_id}، \code{admission\_source\_id} و \code{discharge\_disposition\_id} به برچسب‌های خوانا نگاشت شدند؛ ۰ مقدار نگاشت‌نشده & نگاشت کدها\\
متغیر هدف دودویی \code{early\_readmission}؛ نشانگر \code{death\_or\_hospice} از شناسه‌های ترخیص ۱۱، ۱۳، ۱۴، ۱۹، ۲۰ و ۲۱ (۲,۴۲۳ مراجعه، ۲/۳۸\%)؛ گروه‌بندی \lr{ICD-9} متغیرهای \code{diag\_1} تا \code{diag\_3} در ده دسته‌ی کلی (گردش خون، تنفسی، گوارشی، دیابت، آسیب، اسکلتی‌عضلانی، ادراری‌تناسلی، نئوپلاسم‌ها، سایر و گمشده) & مهندسی ویژگی\\
۵۱ ستون پیش‌بینی‌کننده (۱۱ عددی، ۴۰ طبقه‌ای). نژاد، جنسیت و سن به‌عنوان پیش‌بینی‌کننده نگه داشته شدند و برای ارزیابی زیرگروه‌ها به کار رفتند. از پیش‌بینی‌کننده‌ها کنار گذاشته شدند: شناسه‌ها، نشانگر بیمار تکراری، ستون‌های پیامد، \code{weight} و دو ویژگی ثابت & پیش‌بینی‌کننده‌ها\\

\bottomrule
\end{tabularx}}
\end{table}

شمار مراجعه‌های با مصرف خدمات بسیار زیاد (برای نمونه تا ۷۶ مراجعه‌ی اورژانس و ۴۲ مراجعه‌ی سرپایی در دوره‌ی پیشین برای یک مراجعه) وجود دارد که حفظ شدند، زیرا ممکن است بیمارانِ واقعاً پرمصرف را نمایندگی کنند. افزون بر این، ۳,۶۹۱ مراجعه (۳/۶۳\%) مقصد ترخیص گمشده دارند و ۵,۲۹۱ نوع پذیرش با مقدار \code{NULL} کدگذاری شده است.

گام‌های زیر انجام \emph{نشد}، زیرا فایل‌های پروژه از آن‌ها پشتیبانی نمی‌کنند: نمونه‌برداری مجدد \po{}\lr{SMOTE} یا کم‌نمونه‌گیری)، انتخاب ویژگی رسمی، تبدیل‌های لگاریتمی یا توانی، تجمیع ستون‌های دارویی، پالایش هم‌خطی و جای‌گذاری فراتر از میانه/پرتکرارترین مقدار.

\subsection{تحلیل آماری اکتشافی}

تحلیل اکتشافی بر اندازه‌های اثر و مقایسه‌های توصیفیِ شفاف تأکید دارد.

\subsubsection{تفاوت میانگین استانداردشده}

برای متغیر عددی $X$، تفاوت میانگین استانداردشده میان مراجعه‌های دارای بستری مجدد زودهنگام و سایر مراجعه‌ها چنین است:
\mathblock{\mathrm{SMD}=\frac{\bar X_1-\bar X_0}{\sqrt{\frac{1}{2}(s_1^2+s_0^2)}}.}
این کمیتِ بدون واحد امکان مقایسه‌ی متغیرهای با مقیاس‌های متفاوت را فراهم می‌کند. مقادیر مطلق \lr{SMD} در حدود ۰/۲، ۰/۵ و ۰/۸ در راهنماهای روش‌شناختی عمومی اغلب به‌ترتیب کوچک، متوسط و بزرگ توصیف می‌شوند، هرچند این آستانه‌ها قانون بالینی نیستند.

\subsubsection{ضریب \lr{V} کرامر}

برای متغیرهای طبقه‌ای، ارتباط با پیامد دودویی با ضریب \lr{V} کرامر خلاصه شد:
\mathblock{V=\sqrt{\frac{\chi^2/n}{\min(r-1,k-1)}},}
که در آن $n$ حجم نمونه و $r$ و $k$ به‌ترتیب تعداد سطرها و ستون‌های جدول پیشامد هستند.

\subsubsection{همبستگی و غربالگری داده‌های پرت}

همبستگی پیرسون به‌عنوان ابزار غربالگری توصیفی برای روابط میان متغیرهای عددیِ مصرف خدمات به کار رفت. مقادیر فرین با حصارهای متعارف دامنه‌ی بین‌چارکی \lr{(IQR)} غربال شدند:
\mathblock{\left[Q_1-1.5\,\mathrm{IQR},\;Q_3+1.5\,\mathrm{IQR}\right],\qquad \mathrm{IQR}=Q_3-Q_1.}
مقادیر فرینِ مصرف خدمات در صورتی که می‌توانستند بیماران واقعاً پرمصرف را نمایندگی کنند، حفظ شدند.

\subsubsection{دامنه‌ی استنباط آماری}

دفترچه‌ی تحلیل به‌جای آزمون‌های فرضِ رسمی، به آماره‌های توصیفیِ مبتنی بر اندازه‌ی اثر تکیه دارد؛ \emph{هیچ مقدار \lr{p}ای در هیچ‌جای تحلیل گزارش نشده است و در اینجا نیز چیزی اضافه نشده است.} چون تحلیل اکتشافی است، فرضیه‌ها به‌طور رسمی بیان نشدند و سطح معناداری‌ای اعمال نمی‌شود. آماره‌ی خی‌دو درون تابع \lr{V} کرامر محاسبه می‌شود اما نمایش داده نمی‌شود. با حدود ۱۰۰,۰۰۰ مراجعه، اندازه‌های اثر انتخاب مناسبی‌اند، زیرا حتی تفاوت‌های ناچیز نیز مقدار \lr{p}های بسیار کوچکی به دست می‌دهند؛ اندازه‌های اثر نشان می‌دهند که آیا تفاوت آن‌قدر بزرگ هست که اهمیت داشته باشد.

\begin{table}[H]
\centering
\caption{روش‌های آماری به‌کاررفته و دلیل انتخاب آن‌ها.}
\label{tab:methods}
\small
{\faindent=0pt\relax\begin{tabularx}{\textwidth}{>{\raggedleft\arraybackslash}X>{\raggedleft\arraybackslash}X>{\beginR}r<{\endR}}
\toprule
دلیل انتخاب & هدف & روش\\

\midrule
اندازه‌ی اثرِ بدون واحد که صرف‌نظر از حجم نمونه آگاهی‌بخش است & مقایسه‌ی متغیرهای عددی میان بستری مجدد زودهنگام و سایر مراجعه‌ها & تفاوت میانگین استانداردشده\\
متغیرهای با تعداد سطوح متفاوت را مدیریت می‌کند؛ بین ۰ و ۱ کراندار است & ارتباط یک متغیر طبقه‌ای با پیامد دودویی & ضریب \lr{V} کرامر\\
غربالگری ساده‌ی ارتباط خطی و افزونگی & روابط میان متغیرهای عددی & همبستگی پیرسون\\
ابزار غربالگری است، نه قاعده‌ی حذف خودکار & غربالگری مقادیر فرین & قاعده‌ی $1.5\times$\lr{IQR}\\
شفاف است؛ اندازه‌ی نمونه‌ها نشان داده می‌شود تا گروه‌های کوچک بیش‌ازحد تفسیر نشوند & مقایسه‌ی درصدهای بستری مجدد زودهنگام میان گروه‌ها & نرخ‌های گروهی همراه با شمارش‌ها\\

\bottomrule
\end{tabularx}}
\end{table}

مقدار $\mathrm{SMD}=0.415$ بر اساس قواعد سرانگشتیِ رایج (۰/۲ کوچک، ۰/۵ متوسط) تفاوتی کوچک تا متوسط است؛ قوی‌ترین سیگنال عددی است ولی همچنان هم‌پوشانی زیادی میان دو گروه باقی می‌گذارد. همه‌ی مقادیر \lr{V} کرامر کوچک‌اند (بزرگ‌ترین آن‌ها ۰/۱۲۵ است) و ارتباط‌هایی که از نظر آماری قابل‌تشخیص‌اند لزوماً از نظر عملی مفید نیستند. داده‌ها مشاهده‌ای‌اند و هیچ طرح علّی به کار نرفته است. یک توسعه‌ی طبیعی که در اینجا انجام نشد، آزمون‌های خی‌دو و مان-ویتنی (یا رگرسیون لجستیک) همراه با بازه‌های اطمینان و تصحیح مقایسه‌های چندگانه است.

\subsection{کنترل نشت در سطح بیمار}

شناسه‌ی بیمار \code{patient\_nbr} فقط برای گروه‌بندی به کار رفت و به‌عنوان ویژگی پیش‌بینی‌کننده وارد نشد. یک تقسیم ۸۰/۲۰ با \lr{GroupShuffleSplit} و \code{random\_state=123} نتایج زیر را داد:
\begin{faitemize}
\faitem آموزش: ۷۹,۳۹۵ مراجعه و ۵۵,۹۹۲ بیمار؛
\faitem آزمون: ۱۹,۹۴۸ مراجعه و ۱۳,۹۹۸ بیمار؛
\faitem هم‌پوشانی بیماران: صفر.
\end{faitemize}
این یک پادمان حیاتی است، چون یک بیمار می‌تواند چند مراجعه‌ی بیمارستانی داشته باشد. طرح مدل‌سازی بدین‌ترتیب تعمیم‌پذیری به بیمارانی را برآورد می‌کند که در بخش آموزش حضور نداشته‌اند.

\subsection{مدل‌های یادگیری ماشین}

چهار الگوریتم مقایسه شدند:
\begin{faenum}
\faitem رگرسیون لجستیک به‌عنوان یک مبنای خطی شفاف؛
\faitem درخت تصمیم به‌عنوان مدلی غیرخطی و تفسیرپذیر؛
\faitem جنگل تصادفی به‌عنوان مجموعه‌ای از درخت‌های تصمیم؛
\faitem \lr{XGBoost} به‌عنوان مجموعه‌ی درخت‌های تقویت‌شده‌ی گرادیانی.
\end{faenum}

عدم توازن کلاس‌ها با وزن‌دهی کلاس برای رگرسیون لجستیک و درخت تصمیم، زیرنمونه‌گیری متوازن برای جنگل تصادفی، و نسبت وزن کلاس مثبت برای \lr{XGBoost} مدیریت شد.

\begin{table}[H]
\centering
\caption{پیکربندی مدل‌های مرحله‌ی الف.}
\label{tab:models}
{\faindent=0pt\relax\begin{tabularx}{\textwidth}{>{\raggedleft\arraybackslash}X>{\beginR}r<{\endR}}
\toprule
تنظیمات اصلی & مدل\\

\midrule
\code{max\_iter=2000}؛ وزن‌های کلاس متوازن & رگرسیون لجستیک\\
\code{max\_depth=8}؛ \code{min\_samples\_leaf=50}؛ وزن‌های کلاس متوازن & درخت تصمیم\\
۲۵۰ درخت؛ \code{max\_depth=16}؛ \code{min\_samples\_leaf=10}؛ زیرنمونه‌گیری متوازن & جنگل تصادفی\\
۲۵۰ درخت؛ عمق ۵؛ نرخ یادگیری ۰/۰۵؛ \code{min\_child\_weight=5}؛ زیرنمونه ۰/۸؛ زیرنمونه‌ی ستونی ۰/۸؛ تنظیم‌سازی ۱/۰ & \lr{XGBoost}\\

\bottomrule
\end{tabularx}}
\end{table}

\subsection{بهینه‌سازی آستانه}

آستانه‌ی احتمال پیش‌فرض ۰/۵۰ وقتی هدف بالینی بر حساسیت تأکید دارد، لزوماً مناسب نیست. بنابراین از یک تقسیم اعتبارسنجی درون داده‌ی آموزش برای جست‌وجوی آستانه‌ها از ۰/۰۵ تا ۰/۵۰ با گام ۰/۰۱ استفاده شد. آستانه‌ای که \lr{$F_2$} را بیشینه می‌کرد، برای هر مدل مرحله‌ی الف انتخاب شد. مجموعه‌ی آزمون نگه‌داشته‌شده در طول این انتخاب دست‌نخورده ماند.

جست‌وجوی آستانه یک نقطه‌ی کاری صریح و مبتنی بر مصالحه می‌سازد: پایین‌آوردن آستانه شمار مراجعه‌های علامت‌گذاری‌شده را افزایش می‌دهد و عموماً بازیابی را بالا و دقت را پایین می‌برد. این مصالحه برای تفسیر مدل به‌عنوان ابزار اولویت‌بندی، و نه یک ابزار تشخیصی خودمختار، محوری است.

\subsection{جست‌وجوی ابرپارامتر}

یک مدل \lr{XGBoost} مرحله‌ی دوم از \code{RandomizedSearchCV} با ۱۲ ترکیب نمونه‌برداری‌شده از پارامترها و اعتبارسنجی سه‌بخشیِ \code{StratifiedGroupKFold} گروه‌بندی‌شده بر اساس بیمار استفاده کرد. معیار بهینه‌سازی \lr{$F_2$} بود. بهترین مجموعه‌ی پارامترها چنین بود:
\mathblock{n_{\mathrm{estimators}}=200,\quad \mathrm{max\usc{}depth}=3,\quad \mathrm{learning\ rate}=0.05,}
\mathblock{\mathrm{min\usc{}child\ weight}=5,\quad \mathrm{subsample}=0.8,\quad \mathrm{colsample\usc{}bytree}=1.0,\quad \lambda=1.0.}
بهترین \lr{$F_2$} اعتبارسنجی‌متقاطع برابر ۰/۴۰۶ بود.

\subsection{تحلیل زیرگروه‌ها}

تحلیل با رویکرد انصاف بر بازیابی و \lr{FNR} درون گروه‌های نژاد، جنسیت و سن متمرکز بود. برای هر مدل:
\mathblock{\mathrm{FNR\ gap}=\max_g(\mathrm{FNR}_g)-\min_g(\mathrm{FNR}_g).}
این برآوردها همراه با اندازه‌ی نمونه‌ی زیرگروه‌ها تفسیر شدند. به‌ویژه، \lr{FNR} برابر ۱/۰ در گروهی که تنها یک مورد مثبت دارد، همان قدرت شواهد را ندارد که \lr{FNR} برآوردشده از صدها مورد مثبت دارد.

% ============================================================
\section{یافته‌ها}
% ============================================================

\subsection{توزیع پیامد بستری مجدد}

توزیع پیامد سه‌سطحی به‌شدت نامتوازن است. از ۱۰۱,۷۶۶ مراجعه، ۱۱,۳۵۷ مورد (۱۱/۱۶\%) ظرف ۳۰ روز با بستری مجدد دنبال شدند، ۳۴/۹۳\% پس از ۳۰ روز بستری مجدد شدند و ۵۳/۹۱\% بستری مجدد نشدند.

\fig{0.78}{Distribution_of_Readmission_Outcomes.png}
\captionof{figure}{توزیع سه پیامد بستری مجدد در کل مجموعه‌داده. کلاس بستری مجدد زودهنگام پیامد اقلیت است.}
\label{fig:outcome}

عدم توازن نشان‌داده‌شده در شکل~\ref{fig:outcome} توضیح می‌دهد که چرا صحت را نمی‌توان تنها معیار عملکرد دانست. پرسش بالینی دقیقاً این است که آیا موارد اقلیتِ بستری مجدد زودهنگام را می‌توان با بار هشدار کاذبِ قابل‌قبول شناسایی کرد.

\subsection{توزیع هدف دودویی}

چون مسئله‌ی مدل‌سازی دودویی است، سه دسته‌ی اصلی ادغام شدند: \code{NO} و \gtthirty{} کلاس ۰ (بدون بستری مجدد زودهنگام) و \lessthirty{} کلاس ۱ را می‌سازند. شکل~\ref{fig:binary} توازن کلاس‌ها را در سه مرحله‌ی تحلیل نشان می‌دهد.

\begin{table}[H]
\centering
\caption{نگاشت دسته‌های اصلی بستری مجدد به هدف دودویی.}
\label{tab:binary}
{\faindent=0pt\relax\begin{tabular}{>{\beginR}r<{\endR}>{\beginR}r<{\endR}>{\beginR}r<{\endR}}
\toprule
هدف دودویی & مراجعه‌ها & دسته‌ی اصلی\\

\midrule
۱ (بستری مجدد زودهنگام) & ۱۱,۳۵۷ (۱۱/۱۶\%) & \lessthirty\\
۰ & ۳۵,۵۴۵ (۳۴/۹۳\%) & \gtthirty\\
۰ & ۵۴,۸۶۴ (۵۳/۹۱\%) & \code{NO}\\
۰ & ۹۰,۴۰۹ (۸۸/۸۴\%) & مجموع کلاس ۰\\

\bottomrule
\end{tabular}}
\end{table}

\fig{0.98}{Distribution_of_Binary_Target_-_Early_Readmission.png}
\captionof{figure}{هدف دودویی (بستری مجدد زودهنگام کمتر از ۳۰ روز در برابر سایر موارد) در کل مجموعه‌داده ($n=101{,}766$)، هم‌گروه مدل‌سازی ($n=99{,}343$) و مجموعه‌ی آزمون نگه‌داشته‌شده ($n=19{,}948$). بستری مجدد زودهنگام در هر مرحله کلاس اقلیت است (۱۱/۱۶\%، ۱۱/۳۹\% و ۱۱/۳۳\%).}
\label{fig:binary}

بستری مجدد زودهنگام در هر مرحله پیامد اقلیت است. طبقه‌بندی‌ای که برای هر مراجعه «بدون بستری مجدد زودهنگام» پیش‌بینی کند حدود ۸۹\% مواقع درست خواهد بود، در حالی که از نظر بالینی بی‌فایده است؛ به همین دلیل صحت معیار اصلی نیست. توزیع سه‌دسته‌ای اصلی (شکل~\ref{fig:outcome}) صرفاً توصیفی است و هدف مدل‌سازی نیست.

\subsection{توزیع سنی}

هم‌گروه در سالمندان متمرکز است. بزرگ‌ترین گروه ده‌ساله ۷۰ تا ۸۰ سال با ۲۶,۰۶۸ مراجعه است، پس از آن ۶۰ تا ۷۰ سال با ۲۲,۴۸۳ و ۸۰ تا ۹۰ سال با ۱۷,۱۹۷ مراجعه قرار دارند.

\fig{0.78}{Age_Distribution_of_Hospital_Encounters.png}
\captionof{figure}{توزیع سنیِ مراجعه‌های بیمارستانی در گروه‌های ده‌ساله.}
\label{fig:age}

\begin{table}[H]
\centering
\caption{نرخ بستری مجدد زودهنگام بر حسب گروه سنی.}
\label{tab:age}
{\faindent=0pt\relax\begin{tabular}{>{\beginR}r<{\endR}>{\beginR}r<{\endR}>{\beginR}r<{\endR}>{\beginR}r<{\endR}}
\toprule
نرخ & بستری‌های مجدد زودهنگام & مراجعه‌ها & گروه سنی\\

\midrule
۱/۸۶\% & ۳ & ۱۶۱ & \lr{[0--10)}\\
۵/۷۹\% & ۴۰ & ۶۹۱ & \lr{[10--20)}\\
۱۴/۲۴\% & ۲۳۶ & ۱,۶۵۷ & \lr{[20--30)}\\
۱۱/۲۳\% & ۴۲۴ & ۳,۷۷۵ & \lr{[30--40)}\\
۱۰/۶۰\% & ۱,۰۲۷ & ۹,۶۸۵ & \lr{[40--50)}\\
۹/۶۷\% & ۱,۶۶۸ & ۱۷,۲۵۶ & \lr{[50--60)}\\
۱۱/۱۳\% & ۲,۵۰۲ & ۲۲,۴۸۳ & \lr{[60--70)}\\
۱۱/۷۷\% & ۳,۰۶۹ & ۲۶,۰۶۸ & \lr{[70--80)}\\
۱۲/۰۸\% & ۲,۰۷۸ & ۱۷,۱۹۷ & \lr{[80--90)}\\
۱۱/۱۰\% & ۳۱۰ & ۲,۷۹۳ & \lr{[90--100)}\\

\bottomrule
\end{tabular}}
\end{table}

\fig{0.78}{Early_Readmission_Rate_by_Age_Group.png}
\captionof{figure}{نرخ بستری مجدد زودهنگام بر حسب گروه سنی. بیشتر گروه‌های بزرگسال در حدود ۱۰ تا ۱۲\% قرار دارند؛ اوج ظاهری در سنین ۲۰ تا ۳۰ بر پایه‌ی زیرگروهی نسبتاً کوچک است.}
\label{fig:agerate}

شباهت گسترده‌ی نرخ‌ها در بزرگسالان نشان می‌دهد که سن به‌تنهایی در این مجموعه‌داده جداکننده‌ی توصیفیِ قوی نیست. درصد نسبتاً بالا در گروه ۲۰ تا ۳۰ سال نباید اثر سن تفسیر شود، زیرا مخرج آن بسیار کوچک‌تر از گروه‌های مسن‌تر است.

\subsection{مصرف قبلی خدمات و طول مدت اقامت}

مصرف قبلی خدمات بستری روشن‌ترین تفکیک توصیفی را نشان می‌دهد. میانگین بستری‌های قبلی در مراجعه‌های دارای بستری مجدد زودهنگام ۱/۲۲ بود، در مقایسه با ۰/۸۴ در مراجعه‌های دارای بستری مجدد دیرهنگام و ۰/۳۸ در مراجعه‌های بدون بستری مجدد.

\fig{0.78}{Prior_Inpatient_Visits_by_Readmission_Group.png}
\captionof{figure}{بستری‌های قبلی بر حسب پیامد بستری مجدد. مراجعه‌های دارای بستری مجدد زودهنگام به سمت مصرف قبلی بیشتر خدمات بستری جابه‌جا شده‌اند.}
\label{fig:inpatient}

طول مدت اقامت تفکیک‌کنندگیِ کمتری دارد. میانگین اقامت ۴/۷۷ روز برای بستری مجدد زودهنگام، ۴/۵۰ روز برای بستری مجدد دیرهنگام و ۴/۲۵ روز برای بدون بستری مجدد بود.

\fig{0.78}{Time_in_Hospital_by_Readmission_Group.png}
\captionof{figure}{مدت اقامت در بیمارستان بر حسب پیامد بستری مجدد. با وجود تفاوت اندک میانگین‌ها، توزیع‌ها هم‌پوشانی زیادی دارند.}
\label{fig:los}

\subsection{تفاوت‌های استانداردشده}

مقدار مطلق \lr{SMD} برای \code{number\_inpatient} برابر ۰/۴۱۵ بود که به‌مراتب بزرگ‌تر از سایر متغیرهای عددی است:
\begin{table}[H]
\centering
\caption{مقادیر مطلق تفاوت میانگین استانداردشده میان بستری مجدد زودهنگام و سایر مراجعه‌ها.}
\label{tab:smd}
{\faindent=0pt\relax\begin{tabular}{>{\beginR}r<{\endR}>{\beginR}r<{\endR}}
\toprule
\lr{SMD} مطلق & متغیر\\

\midrule
\textbf{۰/۴۱۵} & \code{number\_inpatient}\\
۰/۱۶۳ & \code{number\_diagnoses}\\
۰/۱۵۷ & \code{number\_emergency}\\
۰/۱۴۰ & \code{time\_in\_hospital}\\
۰/۱۲۲ & \code{num\_medications}\\
۰/۰۶۵ & \code{num\_lab\_procedures}\\
۰/۰۵۹ & \code{number\_outpatient}\\

\bottomrule
\end{tabular}}
\end{table}

\fig{0.78}{Standardized_Differences_-_Early_Readmission_vs_Other_Encounters.png}
\captionof{figure}{مقادیر مطلق تفاوت‌های میانگین استانداردشده. مصرف قبلی خدمات بستری قوی‌ترین جداکننده‌ی عددی است.}
\label{fig:smd}

مقدار ۰/۴۱۵ یک تفاوت استانداردشده‌ی کوچک تا متوسط است. مهم این است که این به معنای ناهم‌پوشانی توزیع‌ها نیست. توزیع تصویری در شکل~\ref{fig:inpatient} هم‌پوشانی قابل‌توجه را تأیید می‌کند.

\subsection{ساختار همبستگی}

متغیرهای عددیِ مصرف خدمات همبستگی‌های دوبه‌دوی ضعیف تا متوسط دارند. بزرگ‌ترین همبستگی پیرسون ۰/۴۷ میان \code{time\_in\_hospital} و \code{num\_medications} است؛ همبستگی میان \code{time\_in\_hospital} و \code{num\_lab\_procedures} برابر ۰/۳۲ و میان بستری‌های قبلی و مراجعه‌های اورژانس قبلی ۰/۲۷ است.

\fig{0.78}{Correlation_Between_Numeric_Utilization_Variables.png}
\captionof{figure}{ماتریس همبستگی پیرسون میان متغیرهای عددی مصرف خدمات. هیچ جفتی همبستگی خطی قوی نشان نمی‌دهد.}
\label{fig:corr}

چون چند متغیر شمارشی با صفرهای فراوان و توزیع چوله‌اند، همبستگی پیرسون باید ابزار غربالگری توصیفی تلقی شود، نه توصیفی کامل از وابستگی.

\subsection{ارتباط‌های طبقه‌ای}

بزرگ‌ترین مقدار \lr{V} کرامر برای مقصد ترخیص مشاهده شد:
\mathblock{V_{\mathrm{discharge}}=0.125.}
مقادیر دیگر بسیار کوچک‌تر بودند.

\begin{table}[H]
\centering
\caption{ضریب \lr{V} کرامر برای متغیرهای طبقه‌ای منتخب و بستری مجدد زودهنگام.}
\label{tab:cramer}
{\faindent=0pt\relax\begin{tabular}{>{\beginR}r<{\endR}>{\beginR}r<{\endR}}
\toprule
\lr{V} کرامر & متغیر\\

\midrule
\textbf{۰/۱۲۵} & \code{discharge\_disposition}\\
۰/۰۴۳ & \code{insulin}\\
۰/۰۳۴ & \code{age}\\
۰/۰۲۷ & \code{diabetesMed}\\
۰/۰۲۳ & \code{admission\_source}\\
۰/۰۱۹ & \code{change}\\
۰/۰۱۹ & \code{A1Cresult}\\
۰/۰۱۶ & \code{admission\_type}\\
۰/۰۱۶ & \code{race}\\
۰/۰۱۲ & \code{max\_glu\_serum}\\
۰/۰۰۳ & \code{gender}\\

\bottomrule
\end{tabular}}
\end{table}

ارتباط مقصد ترخیص برای پیش‌بینی آگاهی‌بخش است اما از نظر بالینی تفسیرش دشوار است، زیرا مقصد بالینی را با دسته‌های اداری درآمیخته است. نباید آن را شاهدی بر این دانست که یک مقصد ترخیصِ مشخص موجب بستری مجدد می‌شود.

\subsection{نژاد و جنسیت}

نرخ‌های خام بستری مجدد زودهنگام برای بزرگ‌ترین گروه‌های نژادی مشابه بود.

\begin{table}[H]
\centering
\caption{نرخ بستری مجدد زودهنگام بر حسب نژاد.}
\label{tab:race}
{\faindent=0pt\relax\begin{tabular}{>{\beginR}r<{\endR}>{\beginR}r<{\endR}>{\beginR}r<{\endR}}
\toprule
نرخ بستری مجدد زودهنگام & مراجعه‌ها & نژاد\\

\midrule
۱۱/۲۹\% & ۷۶,۰۹۹ & قفقازی \lr{(Caucasian)}\\
۱۱/۲۲\% & ۱۹,۲۱۰ & آفریقایی‌آمریکایی \lr{(African American)}\\
۱۰/۴۱\% & ۲,۰۳۷ & اسپانیایی‌تبار \lr{(Hispanic)}\\
۱۰/۱۴\% & ۶۴۱ & آسیایی \lr{(Asian)}\\
۹/۶۳\% & ۱,۵۰۶ & سایر \lr{(Other)}\\
۸/۲۷\% & ۲,۲۷۳ & گمشده \lr{(Missing)}\\

\bottomrule
\end{tabular}}
\end{table}

\fig{0.78}{Early_Readmission_Rate_by_Race.png}
\captionof{figure}{نرخ‌های خام بستری مجدد زودهنگام بر حسب نژاد. گروه‌های کوچک‌تر عدم قطعیتِ نمونه‌برداریِ وسیع‌تری دارند که در نمودار توصیفی نشان داده نشده است.}
\label{fig:race}

نرخ بستری مجدد زودهنگام در مراجعه‌های زنان ۱۱/۲۵\% و در مراجعه‌های مردان ۱۱/۰۶\% بود.

\fig{0.72}{Early_Readmission_Rate_by_Gender.png}
\captionof{figure}{نرخ بستری مجدد زودهنگام بر حسب جنسیت. نرخ‌های خام بسیار شبیه هم‌اند.}
\label{fig:gender}

\subsection{گروه تشخیص اصلی}

در میان گروه‌های تشخیصیِ پرشمار، «دیابت» با ۱۲/۹۸\% (۸,۷۵۷ مراجعه) بالاترین نرخ بستری مجدد زودهنگام را داشت و پس از آن «آسیب» با ۱۲/۲۵\% قرار گرفت. گروه‌های تنفسی و اسکلتی‌عضلانی به‌ترتیب نرخ‌های ۹/۷۳\% و ۹/۵۰\% داشتند.

\fig{0.78}{Early_Readmission_Rate_by_Primary_Diagnosis_Group.png}
\captionof{figure}{نرخ بستری مجدد زودهنگام بر حسب گروه تشخیص اصلی. گروه بسیار کوچکِ «گمشده» نباید تفسیر شود.}
\label{fig:diagnosis}

این نرخ‌های گروهی توصیفی‌اند. دسته‌ی تشخیص خلاصه‌ای از اطلاعات کدگذاری‌شده است و ممکن است شدت بالینی، مرحله‌ی بیماری، وضعیت آزمایشگاهی یا عوامل اجتماعی تعیین‌کننده‌ی سلامت را منعکس نکند.

% ============================================================
\section{نتایج مدل‌سازی پیش‌بینی}
% ============================================================

\subsection{مقایسه‌ی مدل‌های مرحله‌ی الف}

در آستانه‌ی پیش‌فرض ۰/۵۰، چهار مدل مقادیر \lr{ROC--AUC} نسبتاً نزدیکی از ۰/۶۵۹ تا ۰/۶۸۳ نشان دادند.

\begin{table}[H]
\centering
\small\setlength{\tabcolsep}{3.5pt}
\caption{عملکرد مدل‌های مرحله‌ی الف در آستانه‌ی احتمال پیش‌فرض ۰/۵۰.}
\label{tab:default}
{\faindent=0pt\relax\begin{tabular}{>{\beginR}r<{\endR}>{\beginR}r<{\endR}>{\beginR}r<{\endR}>{\beginR}r<{\endR}>{\beginR}r<{\endR}>{\beginR}r<{\endR}>{\beginR}r<{\endR}>{\beginR}r<{\endR}}
\toprule
\lr{PR--AUC} & \lr{ROC--AUC} & \lr{F1} & \lr{Precision} & \lr{Recall} & \lr{Bal. Acc.} & \lr{Accuracy} & مدل\\

\midrule
۰/۲۲۱ & ۰/۶۶۴ & ۰/۲۷۲ & ۰/۱۷۸ & ۰/۵۶۸ & ۰/۶۱۷ & ۰/۶۵۵ & رگرسیون لجستیک\\
۰/۲۱۲ & ۰/۶۵۹ & ۰/۲۶۷ & ۰/۱۶۷ & ۰/۶۵۹ & ۰/۶۲۰ & ۰/۵۹۰ & درخت تصمیم\\
۰/۲۰۷ & ۰/۶۷۴ & ۰/۲۷۷ & ۰/۱۸۰ & ۰/۶۰۳ & ۰/۶۲۶ & ۰/۶۴۳ & جنگل تصادفی\\
\textbf{۰/۲۳۸} & \textbf{۰/۶۸۳} & ۰/۲۸۷ & ۰/۱۸۸ & ۰/۵۹۸ & ۰/۶۳۴ & ۰/۶۶۳ & \lr{XGBoost}\\

\bottomrule
\end{tabular}}
\end{table}

\lr{XGBoost} در آستانه‌ی پیش‌فرض بهترین \lr{ROC--AUC}، \lr{PR--AUC}، صحت متوازن و \lr{F1} را دارد. با این حال تفاوت‌ها اندک‌اند و نتایج را نباید شاهدی بر برتری بزرگ یک الگوریتم نسبت به دیگری دانست.

\subsection{نقاط کاری بازیابی‌محور}

تنظیم آستانه یک نقطه‌ی کاری بالینیِ متفاوت پدید می‌آورد.

\begin{table}[H]
\centering
\small\setlength{\tabcolsep}{3pt}
\caption{عملکرد مرحله‌ی الف در آستانه‌های انتخاب‌شده با \lr{$F_2$} روی مجموعه‌ی اعتبارسنجی.}
\label{tab:tuned}
{\faindent=0pt\relax\begin{tabular}{>{\beginR}r<{\endR}>{\beginR}r<{\endR}>{\beginR}r<{\endR}>{\beginR}r<{\endR}>{\beginR}r<{\endR}>{\beginR}r<{\endR}>{\beginR}r<{\endR}>{\beginR}r<{\endR}>{\beginR}r<{\endR}}
\toprule
\lr{PR--AUC} & \lr{ROC--AUC} & \lr{F1} & \lr{Precision} & \lr{Recall} & \lr{Bal. Acc.} & \lr{Acc.} & \lr{Thr.} & مدل\\

\midrule
۰/۲۲۱ & ۰/۶۶۴ & ۰/۲۴۵ & ۰/۱۴۵ & ۰/۷۹۱ & ۰/۵۹۷ & ۰/۴۴۸ & ۰/۳۹ & رگرسیون لجستیک\\
۰/۲۱۲ & ۰/۶۵۹ & ۰/۲۴۳ & ۰/۱۴۳ & ۰/۷۸۶ & ۰/۵۹۳ & ۰/۴۴۴ & ۰/۳۸ & درخت تصمیم\\
۰/۲۰۷ & ۰/۶۷۴ & ۰/۲۴۷ & ۰/۱۴۵ & \textbf{۰/۸۳۰} & ۰/۶۰۲ & ۰/۴۲۶ & ۰/۴۵ & جنگل تصادفی\\
\textbf{۰/۲۳۸} & \textbf{۰/۶۸۳} & \textbf{۰/۲۵۵} & \textbf{۰/۱۵۲} & ۰/۷۹۷ & \textbf{۰/۶۱۵} & ۰/۴۷۴ & ۰/۴۱ & \lr{XGBoost}\\

\bottomrule
\end{tabular}}
\end{table}

این نتیجه یک مصالحه‌ی عملیاتی محوری را نشان می‌دهد. برای \lr{XGBoost}، حرکت از آستانه‌ی پیش‌فرض به ۰/۴۱ بازیابی را از ۰/۵۹۸ به ۰/۷۹۷ می‌رساند، اما دقت تنها ۰/۱۵۲ می‌ماند و صحت به ۰/۴۷۴ می‌افتد. بدین‌ترتیب در این نقطه‌ی کاری تقریباً شش مراجعه از هر ده مراجعه‌ی آزمون علامت‌گذاری می‌شود.

\subsection{آستانه‌ی تصمیم: اثر آن بر \texorpdfstring{\lr{XGBoost}}{XGBoost}}

یک طبقه‌بند احتمالاتی امتیاز ریسک تولید می‌کند؛ آستانه‌ی تصمیم تعیین می‌کند که از چه امتیازی به بالا، مراجعه به‌عنوان بستری مجدد زودهنگامِ محتمل علامت‌گذاری شود. با شیوع پیامد در حدود ۱۱\%، حد جداکننده‌ی پیش‌فرض ۰/۵۰ نقطه‌ی کاری نامناسبی است، وقتی از دست دادن یک بستری مجدد زودهنگام واقعی گران‌تر از یک هشدار کاذب باشد. بنابراین آستانه به‌صراحت انتخاب شد: برای هر مدل، مقداری در بازه‌ی $[0.05,\,0.50]$ (با گام ۰/۰۱) که \lr{$F_2$} را روی یک مجموعه‌ی اعتبارسنجیِ جداشده بر اساس بیمار بیشینه می‌کرد برگزیده شد، مدل روی کل مجموعه‌ی آموزش دوباره برازش شد، و مجموعه‌ی آزمون هرگز برای انتخاب آستانه به کار نرفت (بدون نشت). برای \lr{XGBoost} آستانه‌ی انتخاب‌شده ۰/۴۱ بود. شکل~\ref{fig:threshold} و جدول~\ref{tab:threshold} مدل \lr{XGBoost} مرحله‌ی الف را روی مجموعه‌ی آزمون نگه‌داشته‌شده، با و بدون این آستانه مقایسه می‌کنند.

\fig{0.95}{XGBoost_-_With_vs_Without_Threshold.png}
\captionof{figure}{\lr{XGBoost} مرحله‌ی الف روی مجموعه‌ی آزمون نگه‌داشته‌شده: آستانه‌ی ۰/۴۱ تنظیم‌شده روی اعتبارسنجی در برابر آستانه‌ی پیش‌فرض ۰/۵۰. \lr{ROC--AUC} و \lr{PR--AUC} به آستانه وابسته نیستند؛ آستانه بازیابی را از ۰/۵۹۸ به ۰/۷۹۷ می‌رساند، به بهای کاهش دقت و صحت.}
\label{fig:threshold}

\begin{table}[H]
\centering
\caption{اثر آستانه‌ی تصمیم بر \lr{XGBoost} مرحله‌ی الف (مجموعه‌ی آزمون نگه‌داشته‌شده). مقادیر علامت‌دار با $^{\dagger}$ از بازیابی، دقت و اندازه‌ی مجموعه‌ی آزمونِ گزارش‌شده استخراج شده‌اند.}
\label{tab:threshold}
\small
{\faindent=0pt\relax\begin{tabular}{>{\beginR}r<{\endR}>{\beginR}r<{\endR}>{\beginR}r<{\endR}>{\raggedleft\arraybackslash}p{5.6cm}}
\toprule
تغییر & پیش‌فرض ۰/۵۰ & آستانه‌ی ۰/۴۱ & معیار\\

\midrule
‎+۰/۱۹۹ (‎+۳۳\% نسبی) & ۰/۵۹۸ & \textbf{۰/۷۹۷} & \lr{Recall} \po{}\lessthirty{}\pc{}\\
‎-۰/۰۳۶ & ۰/۱۸۸ & ۰/۱۵۲ & \lr{Precision} \po{}\lessthirty{}\pc{}\\
‎-۰/۰۳۲ & ۰/۲۸۷ & ۰/۲۵۵ & \lr{F1} \po{}\lessthirty{}\pc{}\\
‎-۰/۰۱۹ & ۰/۶۳۴ & ۰/۶۱۵ & صحت متوازن\\
‎-۰/۱۸۹ & ۰/۶۶۳ & ۰/۴۷۴ & صحت \lr{(Accuracy)}\\
۰ (مستقل از آستانه) & ۰/۶۸۳ & ۰/۶۸۳ & \lr{ROC--AUC}\\
۰ (مستقل از آستانه) & ۰/۲۳۸ & ۰/۲۳۸ & \lr{PR--AUC}\\
‎+۰/۰۱۵ & ۰/۴۱۶ & ۰/۴۳۱ & \lr{$F_2$} (استخراج‌شده)$^{\dagger}$\\
\ensuremath{\approx}‎+۴۵۱ & \ensuremath{\approx}۱,۳۵۱$^{\dagger}$ & ۱,۸۰۲ & بستری‌های مجدد زودهنگام واقعیِ شناسایی‌شده (از ۲,۲۶۰)\\
\ensuremath{\approx}‎+۴,۶۵۰ & \ensuremath{\approx}۷,۲۰۰ (\ensuremath{\approx}۳۶\%)$^{\dagger}$ & ۱۱,۸۴۶ (۵۹/۴\%) & مراجعه‌های علامت‌گذاری‌شده (از ۱۹,۹۴۸)\\

\bottomrule
\end{tabular}}
\end{table}

مقادیر استخراج‌شده خطای گرد‌کردنی در حدود $\pm0.5\%$ پایه دارند. حرکت از ۰/۵۰ به ۰/۴۱ نحوه‌ی \emph{رتبه‌بندی} بیماران توسط مدل را تغییر نمی‌دهد \po{}\lr{ROC--AUC} و \lr{PR--AUC} یکسان‌اند)؛ جایی را که هشدار در آن فعال می‌شود تغییر می‌دهد. این آستانه حدود ۴۵۰ بستری مجدد زودهنگام واقعیِ بیشتر را بازیابی می‌کند (بازیابی از ۰/۵۹۸ به ۰/۷۹۷) و بهایش حدود ده هشدار اضافه به ازای هر مورد اضافیِ شناسایی‌شده است (استخراج‌شده). این تنها انتخاب طراحی‌ای است که بیش از همه توان مدل را در شناسایی موارد مهم بهبود می‌دهد، و همان است که صحت را پایین می‌آورد؛ به همین دلیل از صحت برای قضاوت درباره‌ی مدل استفاده نمی‌شود. آستانه‌ی مناسب به ظرفیت بازبینی بستگی دارد: بیمارستانی که می‌تواند هشدارهای کمتری را بازبینی کند باید آستانه‌ی بالاتری برگزیند.

\subsection{ماتریس‌های درهم‌ریختگی}

ماتریس‌های درهم‌ریختگی مصالحه را ملموس می‌کنند.

\fig{0.82}{Confusion_Matrices_-_All_Models.png}
\captionof{figure}{ماتریس‌های درهم‌ریختگی در آستانه‌های بازیابی‌محور. بازیابی بالا به بهای شمار زیادی هشدار مثبت کاذب به دست می‌آید.}
\label{fig:cm}

برای \lr{XGBoost} تنظیم‌شده، مجموعه‌ی آزمون ۱,۸۰۲ مثبت درست، ۴۵۸ منفی کاذب، ۱۰,۰۴۴ مثبت کاذب و ۷,۶۴۴ منفی درست دارد. \lr{FNR} حاصل برابر است با:
\mathblock{\mathrm{FNR}=\frac{458}{458+1802}=0.203.}
شمار زیاد مثبت‌های کاذب به این معناست که پیاده‌سازی بالینی ظرفیت بازبینی قابل‌توجهی می‌طلبد.

\subsection{اهمیت ویژگی‌ها}

در مدل‌های مبتنی بر درخت، سابقه‌ی بستری قبلی و مقصد ترخیص به‌طور مداوم در میان قوی‌ترین پیش‌بینی‌کننده‌ها دیده می‌شوند.

\begin{table}[H]
\centering
\caption{پیش‌بینی‌کننده‌های پیشرو بر حسب مدل.}
\label{tab:features}
{\faindent=0pt\relax\begin{tabularx}{\textwidth}{>{\raggedleft\arraybackslash}X>{\beginR}r<{\endR}}
\toprule
قوی‌ترین ویژگی‌های گزارش‌شده & مدل\\

\midrule
\code{number\_inpatient} (۰/۴۲۵)؛ مقصد ترخیص «خانه» (۰/۱۳۵)؛ «توان‌بخشی» (۰/۰۵۵) & درخت تصمیم\\
\code{number\_inpatient} (۰/۱۳۲)؛ \code{discharge\_disposition\_id} (۰/۰۵۴)؛ «خانه» (۰/۰۵۱) & جنگل تصادفی\\
«خانه» (۰/۰۲۰)؛ \code{number\_inpatient} (۰/۰۱۹)؛ «توان‌بخشی» (۰/۰۰۹) & \lr{XGBoost}\\
ضرایب کدهای تشخیصیِ نادر، از جمله \code{diag\_2\_136} و \code{diag\_3\_537} & رگرسیون لجستیک\\

\bottomrule
\end{tabularx}}
\end{table}

\fig{0.88}{Top_Features_-_All_Models.png}
\captionof{figure}{برترین ویژگی‌ها و ضرایب در چهار مدل مرحله‌ی الف. اهمیت نشان‌دهنده‌ی سهم در مدل است، نه اثر علّی.}
\label{fig:features}

نتیجه‌ی رگرسیون لجستیک نیازمند احتیاط است، زیرا ماتریس یک‌داغِ بسیار پهن شامل سطوح تشخیصیِ نادر است. ضرایب بزرگ برای دسته‌های نادر می‌توانند ناپایدار باشند و بدون تکرار‌پذیری نباید به‌عنوان عوامل خطر بالینیِ استوار تفسیر شوند.

\subsection{تحلیل خطای زیرگروه‌ها}

\begin{table}[H]
\centering
\caption{شکاف \lr{FNR} زیرگروه‌ها در مرحله‌ی الف در آستانه‌های بازیابی‌محور.}
\label{tab:fairness}
{\faindent=0pt\relax\begin{tabular}{>{\beginR}r<{\endR}>{\beginR}r<{\endR}>{\beginR}r<{\endR}>{\beginR}r<{\endR}}
\toprule
شکاف \lr{FNR} جنسیت & شکاف \lr{FNR} سن & شکاف \lr{FNR} نژاد & مدل\\

\midrule
۰/۰۰۶ & ۰/۹۱۰ & ۰/۲۸۸ & رگرسیون لجستیک\\
۰/۰۳۱ & ۰/۹۳۳ & ۰/۰۸۳ & درخت تصمیم\\
۰/۰۲۴ & ۰/۹۶۷ & ۰/۱۱۸ & جنگل تصادفی\\
۰/۰۲۹ & ۰/۸۸۵ & ۰/۱۱۲ & \lr{XGBoost}\\

\bottomrule
\end{tabular}}
\end{table}

\fig{0.84}{FNR_Gap_Across_Patient_Subgroups.png}
\captionof{figure}{تفاوت بیشینه و کمینه‌ی \lr{FNR} درون زیرگروه‌های نژاد، جنسیت و سن.}
\label{fig:fairness}

شکاف‌های گروه‌های سنی از نظر عددی بزرگ‌اند، اما شمار موارد مثبت زیربنایی اندک است. برای مثال، مجموعه‌ی آزمونِ بازتولیدشده در گروه \lr{[0--10)} فقط ۲۷ مراجعه و یک مورد بستری مجدد زودهنگام، و در گروه \lr{[10--20)} ۱۱۹ مراجعه با شش مورد مثبت دارد. بنابراین \lr{FNR} برابر ۱/۰ می‌تواند از از دست رفتن تنها یک رویداد ناشی شود. به همین ترتیب، دسته‌ی «گمشده» در جنسیت فقط یک مراجعه و هیچ مورد مثبت ندارد و \lr{FNR} آن از نظر بالینیِ معنادار تعریف‌نشده است. این یافته‌ها اهمیت گزارش اندازه‌ی نمونه‌ی زیرگروه‌ها در کنار معیارهای نابرابری را نشان می‌دهند.

\subsubsection{نابرابری گروه‌های سنی و پشتوانه‌ی آماری آن}

نتایج گروه‌های سنی جدی‌ترین یافته‌ی انصاف در این پروژه است. هر چهار مدل دست‌کم در یک گروه سنی همه‌ی بستری‌های مجدد زودهنگام را از دست می‌دهند \lr{(FNR = 1.00)} و شکاف \lr{FNR} میان گروه‌های سنی ۰/۸۸۵ تا ۰/۹۶۷ است. جدول~\ref{tab:agepos} شمار مراجعه‌های آزمون و بستری‌های مجدد زودهنگام واقعی را در هر گروه نشان می‌دهد.

\begin{table}[H]
\centering
\caption{مجموعه‌ی آزمون نگه‌داشته‌شده: مراجعه‌ها و بستری‌های مجدد زودهنگام واقعی (موارد مثبت) در هر گروه سنی. با اجرای دوباره‌ی تقسیم گروهی با \code{random\_state=123} استخراج شده است؛ مجموع‌ها با دفترچه‌ی تحلیل مطابقت دارند.}
\label{tab:agepos}
{\faindent=0pt\relax\begin{tabular}{>{\beginR}r<{\endR}>{\beginR}r<{\endR}>{\beginR}r<{\endR}}
\toprule
بستری‌های مجدد زودهنگام واقعی & مراجعه‌های آزمون & گروه سنی\\

\midrule
\textbf{۱} & ۲۷ & \lr{[0--10)}\\
\textbf{۶} & ۱۱۹ & \lr{[10--20)}\\
۷۸ & ۳۹۶ & \lr{[20--30)}\\
۸۳ & ۷۴۳ & \lr{[30--40)}\\
۱۷۲ & ۱,۸۳۷ & \lr{[40--50)}\\
۲۹۸ & ۳,۳۸۷ & \lr{[50--60)}\\
۵۱۳ & ۴,۵۰۴ & \lr{[60--70)}\\
۶۱۱ & ۵,۱۲۷ & \lr{[70--80)}\\
۴۳۸ & ۳,۲۸۱ & \lr{[80--90)}\\
۶۰ & ۵۲۷ & \lr{[90--100)}\\

\midrule
۲,۲۶۰ & ۱۹,۹۴۸ & مجموع\\

\bottomrule
\end{tabular}}
\end{table}

\fig{0.98}{Age_Subgroup_Disparity_and_Support.png}
\captionof{figure}{نابرابری زیرگروه‌های سنی و پشتوانه‌ی آماری آن. چپ: شمار بستری‌های مجدد زودهنگام واقعی در هر گروه سنی در مجموعه‌ی آزمون (قرمز = کمتر از ۱۰). راست: شکاف \lr{FNR} (بیشینه منهای کمینه) برای نژاد، جنسیت و سن در \lr{XGBoost} مرحله‌ی الف با آستانه‌ی ۰/۴۱.}
\label{fig:agesupport}

با تنها ۱ و ۶ مورد مثبت در دو گروه سنیِ جوان‌تر، \lr{FNR} برابر ۱/۰۰ می‌تواند از یک مورد از دست رفته ناشی شود؛ پس \emph{بزرگی دقیق} شکاف نامطمئن است و بازه‌های اطمینان در سطح زیرگروه محاسبه نشده‌اند. این عدم قطعیت اندازه‌ی برآورد را مشروط می‌کند، اما یافته را از بین نمی‌برد. مدل هیچ شاهدی بر توانایی شناسایی بستری‌های مجدد زودهنگام در بیماران زیر ۲۰ سال نشان نمی‌دهد، و شکل خودِ دفترچه‌ی تحلیل (شکل~\ref{fig:fairness}) نیز آن را نابرابری بحرانی برچسب می‌زند. پیامدهای عملی عبارت‌اند از:
\begin{faenum}
\faitem مدل نباید بدون اعتبارسنجی جداگانه برای بیماران زیر ۲۰ سال به کار رود؛
\faitem بازیابی در گروه‌های سنی باید با بازه‌های اطمینان گزارش شود؛
\faitem کالیبراسیون آگاه از سن یا آستانه‌های ویژه‌ی هر گروه سنی باید در نظر گرفته شود؛
\faitem باید موارد بیشتری در این گروه‌ها گردآوری یا ادغام شود.
\end{faenum}
شکاف‌های نژاد متوسط‌اند (بزرگ‌ترین آن‌ها برای رگرسیون لجستیک، ۰/۲۸۸) و شکاف‌های جنسیت کوچک‌اند (۰/۰۰۶ تا ۰/۰۳۱). مقدار $\min\usc{}\mathrm{Recall}=0.000$ گزارش‌شده برای جنسیت از گروه «گمشده»ی جنسیت با یک مراجعه و بدون مورد مثبت ناشی می‌شود (که با \code{zero\_division=0} امتیاز صفر گرفته است) و شکست مدل نیست.

\subsection{مدل نهایی \texorpdfstring{\lr{XGBoost}}{XGBoost} در مرحله‌ی ب}

مدل \lr{XGBoost} مرحله‌ی دوم با اعتبارسنجی متقاطع تصادفیِ گروهی بهینه شد. بهترین ترکیب پارامترهای آن \lr{$F_2$} اعتبارسنجی‌متقاطع برابر ۰/۴۰۶ داد. هنگامی که در آستانه‌ی ثابت ۰/۵۰ روی مجموعه‌ی آزمون نگه‌داشته‌شده ارزیابی شد، نتایج زیر به دست آمد:

\begin{table}[H]
\centering
\caption{مدل نهایی \lr{XGBoost} با ابرپارامترهای جست‌وجوشده روی مجموعه‌ی آزمون نگه‌داشته‌شده.}
\label{tab:finalxgb}
{\faindent=0pt\relax\begin{tabular}{>{\beginR}r<{\endR}>{\beginR}r<{\endR}}
\toprule
مقدار & معیار\\

\midrule
۰/۶۴۰ & صحت \lr{(Accuracy)}\\
۰/۶۱۹ & \lr{Recall}\\
۰/۱۸۱ & \lr{Precision}\\
۰/۲۸۰ & \lr{F1}\\
$\approx 0.68$ & \lr{ROC--AUC}\\
۰/۵۵۵ & میانگین کلان دقت\\
۰/۶۳۱ & میانگین کلان بازیابی\\
۰/۵۲۰ & میانگین کلان \lr{F1}\\

\bottomrule
\end{tabular}}
\end{table}

نتیجه‌ی مرحله‌ی ب را نباید با \lr{XGBoost} مرحله‌ی الف در آستانه‌ی ۰/۴۱ اشتباه گرفت. این‌ها پیکربندی‌ها و نقاط کاریِ متفاوتی هستند. مرحله‌ی الف برای نشان دادن مصالحه‌ی بازیابی-دقت مفید است؛ مرحله‌ی ب پیکربندی نهاییِ جست‌وجوشده‌ای است که پروژه گزارش کرده است.

\fig{0.60}{ROC_Curve_-_XGBoost.png}
\captionof{figure}{منحنی \lr{ROC} برای مدل نهایی \lr{XGBoost}. توان تفکیک گزارش‌شده متوسط است و \lr{ROC--AUC} حدود ۰/۶۸ است.}
\label{fig:roc}

% ============================================================
\section{یافته‌های کلیدی}
% ============================================================

\begin{faenum}
\faitem \textbf{مسئله دودویی و نامتوازن است:} ۱۱/۱۶\% از ۱۰۱,۷۶۶ مراجعه ظرف ۳۰ روز با بستری مجدد دنبال شدند (۱۱/۳۹\% در هم‌گروه مدل‌سازی با ۹۹,۳۴۳ مراجعه؛ ۱۱/۳۳\% در مجموعه‌ی آزمون).
\faitem \textbf{آستانه‌ی تصمیم، نقطه‌ی قوت کلیدی مدل است:} \lr{XGBoost} در آستانه‌ی ۰/۴۱ تنظیم‌شده روی اعتبارسنجی به بازیابی ۰/۷۹۷ در برابر ۰/۵۹۸ در آستانه‌ی پیش‌فرض ۰/۵۰ می‌رسد، با \lr{ROC--AUC} (۰/۶۸۳) و \lr{PR--AUC} (۰/۲۳۸) یکسان. بهای آن دقت ۰/۱۵۲ (در برابر ۰/۱۸۸) و صحت ۰/۴۷۴ (در برابر ۰/۶۶۳) است.
\faitem \textbf{نابرابری شدید در گروه‌های سنی:} شکاف \lr{FNR} میان گروه‌های سنی در هر چهار مدل ۰/۸۸۵ تا ۰/۹۶۷ است و دست‌کم یک گروه \lr{FNR = 1.00} دارد. این بزرگ‌ترین مسئله‌ی انصاف در پروژه است؛ گروه‌های درگیر فقط ۱ و ۶ مورد مثبت دارند، بنابراین اندازه‌ی دقیق آن نامطمئن است.
\faitem \textbf{سابقه‌ی بستری قبلی روشن‌ترین سیگنال است:} \code{number\_inpatient} بزرگ‌ترین \lr{SMD} (۰/۴۱۵) را دارد؛ میانگین بستری‌های قبلی ۱/۲۲ برای \lessthirty{} در برابر ۰/۳۸ برای \code{NO} است. همچنین در همه‌ی مدل‌های مبتنی بر درخت ویژگی اول یا دوم است.
\faitem \textbf{بیشتر ارتباط‌های دیگر ضعیف‌اند:} همه‌ی مقادیر \lr{V} کرامر $\le 0.125$ هستند؛ ارتباط‌های گروه‌های تشخیص حدود ۰/۰۲ تا ۰/۰۳ است.
\faitem \textbf{مقصد ترخیص آگاهی‌بخش اما دشوارتفسیر است} ($V=0.125$؛ در همه‌ی مدل‌های درختی برجسته است)، چون مقصد بالینی را با دسته‌های اداری درمی‌آمیزد.
\faitem \textbf{صحت گمراه‌کننده است:} بالا بردن بازیابی تا حدود ۰/۸۰ صحت را به حدود ۰/۴۳ تا ۰/۴۷ می‌رساند و دقت را در حدود ۰/۱۴ تا ۰/۱۵ می‌گذارد.
\faitem \textbf{مدل‌ها از نظر توان تفکیک به هم نزدیک‌اند:} \lr{ROC--AUC} از ۰/۶۵۹ تا ۰/۶۸۳؛ \lr{XGBoost} در \lr{ROC--AUC} (۰/۶۸۳) و \lr{PR--AUC} (۰/۲۳۸) بهترین است.
\faitem \textbf{یک \lr{XGBoost} با ابرپارامترهای جست‌وجوشده} (تکمیلی، آستانه‌ی ۰/۵۰) به بازیابی ۰/۶۱۹، دقت ۰/۱۸۱، \lr{F1} برابر ۰/۲۸۰ و \lr{ROC--AUC} حدود ۰/۶۸ می‌رسد.
\faitem \textbf{شکاف‌های نژاد و جنسیت به‌ترتیب متوسط و کوچک‌اند} (شکاف \lr{FNR}: ۰/۰۸۳ تا ۰/۲۸۸ و ۰/۰۰۶ تا ۰/۰۳۱).
\faitem \textbf{نشت در سطح بیمار کنترل شد:} هیچ بیماری هم در آموزش و هم در آزمون حضور ندارد، اعتبارسنجی متقاطع گروهی بر اساس بیمار بود و آستانه بدون استفاده از مجموعه‌ی آزمون انتخاب شد.
\end{faenum}

% ============================================================
\section{تفسیر بالینی}
% ============================================================

\subsection{مصرف قبلی خدمات بستری}

پایدارترین سیگنال، مصرف قبلی خدمات بستری است. تحلیل اندازه‌ی اثر، میانگین‌های گروهی و اهمیت ویژگی در مدل‌های درختی همگی در یک جهت‌اند. این از نظر بالینی به‌عنوان نشانگر نیازهای درمانیِ پیچیده یا ناپایدار موجه است، اما داده‌های مشاهده‌ای نمی‌توانند ثابت کنند که خودِ بستری قبلی موجب بستری مجدد آینده می‌شود.

برای کاربرد عملیاتی، مصرف قبلی خدمات بستری را می‌توان نشانگر طبقه‌بندی ریسک در نظر گرفت. تیم مدیریت مراقبت می‌تواند آن را یکی از اجزای یک فرایند بازبینیِ گسترده‌تر به کار گیرد، به‌ویژه همراه با برنامه‌ریزی ترخیص و سایر اطلاعات بالینی که در این مجموعه‌داده به‌طور کامل نمایندگی نشده‌اند.

\subsection{مقصد ترخیص}

مقصد ترخیص بزرگ‌ترین ارتباط طبقه‌ای را دارد و در همه‌ی مدل‌های درختی برجسته است. با این حال باید محتاطانه تفسیر شود. دسته‌های مقصد می‌توانند نقطه‌ی پایانی اداریِ یک بستری را کدگذاری کنند و ممکن است اطلاعاتی را بازتاب دهند که بسیار دیر در مسیر مراقبت در دسترس می‌شود. همچنین می‌توانند مقصد بالینی را با ساختار کدگذاری درآمیزند. در نتیجه، اهمیت ویژگی به این معنا نیست که تغییر مقصد ترخیص لزوماً ریسک بستری مجدد را تغییر دهد.

\subsection{معنای عملی دقت}

دقت ۰/۱۵۲ در آستانه‌ی بازیابی‌محور \lr{XGBoost} یعنی در مجموعه‌ی آزمون نگه‌داشته‌شده تنها حدود ۱۵ مورد از هر ۱۰۰ مراجعه‌ی علامت‌گذاری‌شده بستری مجدد زودهنگام‌اند. بقیه‌ی هشدارها مثبت کاذب‌اند. این در یک سامانه‌ی غربالگری لزوماً غیرقابل‌قبول نیست، اما حجم کار لازم از پرستاران، هماهنگ‌کنندگان مراقبت یا پزشکان را تعیین می‌کند.

در آستانه‌ی ۰/۴۱ \lr{XGBoost} مرحله‌ی الف، ۱۱,۸۴۶ مراجعه از ۱۹,۹۴۸ مراجعه‌ی آزمون علامت‌گذاری شد، یعنی ۵۹/۴\% مجموعه‌ی آزمون. این نقطه‌ی کاری ۷۹/۷\% از بستری‌های مجدد زودهنگام واقعی را می‌گیرد اما صف بازبینی بسیار بزرگی می‌سازد. بنابراین استقرار واقعی به یک قید صریح ظرفیت و ارزیابی آینده‌نگر ارزشِ هر هشدار نیاز دارد.

\textbf{نسبت افزایش \lr{(lift)} نسبت به نرخ پایه (استخراج‌شده).} شیوع در مجموعه‌ی آزمون ۱۱/۳۳\% است. دقت ۰/۱۵۲ \lr{XGBoost} مرحله‌ی الف در آستانه‌ی ۰/۴۱ حدود ۱/۳ برابر نرخ پایه است، در حالی که ۰/۱۸۸ در آستانه‌ی پیش‌فرض ۰/۵۰ حدود ۱/۷ برابر و ۰/۱۸۱ مدل مرحله‌ی ب حدود ۱/۶ برابر است. بنابراین آستانه بازیابی می‌خرد، نه دقت: سیگنال آماری واقعی است اما از نظر عملیاتی پرهزینه است. در آستانه‌ی پیش‌فرض، مدل تقریباً ۳۶\% مراجعه‌ها را علامت می‌زند (استخراج‌شده) و ۵۹/۸\% بستری‌های مجدد زودهنگام را می‌گیرد، در مقابل ۵۹/۴\% علامت‌گذاری‌شده و ۷۹/۷\% گرفته‌شده در ۰/۴۱.

\subsection{کاربرد مدل به‌عنوان پشتیبان تصمیم}

مدل را باید ابزار رتبه‌بندی یا اولویت‌بندی دانست. بیماری را تشخیص نمی‌دهد، درمان تجویز نمی‌کند و ثابت نمی‌کند که مداخله‌ای از بستری مجدد جلوگیری خواهد کرد. بازبینی انسانی ضروری است، به‌ویژه چون مجموعه‌داده تاریخی، اداری و از نظر شدت بالینی و زمینه‌ی اجتماعی ناقص است.

% ============================================================
\section{انصاف، قابلیت اتکا و عدم قطعیت}
% ============================================================

ارزیابی الگوریتمی در حوزه‌ی سلامت به چیزی فراتر از یک معیار کلی نیاز دارد. مدلی ممکن است توان تفکیک کلی قابل‌قبولی داشته باشد و در عین حال موارد را به‌طور نامتناسب در یک زیرگروه از دست بدهد. بنابراین پروژه‌ی حاضر شکاف‌های \lr{FNR} را برای نژاد، جنسیت و سن گزارش می‌کند.

شکاف‌های نژاد در مدل‌های مرحله‌ی الف از ۰/۰۸۳ تا ۰/۲۸۸ و شکاف‌های جنسیت از ۰/۰۰۶ تا ۰/۰۳۱ بود. شکاف‌های سنی بسیار بزرگ‌تر بود، از ۰/۸۸۵ تا ۰/۹۶۷. اما نتایج سنی تحت تأثیر شمار اندک موارد مثبت‌اند. این تمایز مهم است: نابرابریِ مشاهده‌شده‌ی بزرگ معادل شواهد قوی از تبعیض نظام‌مند نیست. با این حال یافته نادیده گرفته نمی‌شود: چون هر مدل دست‌کم در یک گروه سنی همه‌ی بستری‌های مجدد زودهنگام را از دست می‌دهد، مدل‌ها نباید بدون اعتبارسنجی جداگانه برای بیماران درگیر استفاده شوند.

چند شکل از عدم قطعیت کمّی نشده باقی مانده است. پروژه از یک تقسیم گروهیِ آموزش/آزمون استفاده کرد، بنابراین بازه‌های اطمینان برای معیارهای عملکرد محاسبه نشد. تقسیم‌های گروهیِ مکرر یا بازه‌های اطمینان بوت‌استرپ توصیف بهتری از تغییرپذیری نمونه‌گیری می‌دهند. منحنی‌های کالیبراسیون و خطای کالیبراسیون مورد انتظار نیز مفیدند، زیرا \lr{ROC--AUC} رتبه‌بندی را می‌سنجد، نه اینکه احتمال‌های پیش‌بینی‌شده با ریسک‌های مشاهده‌شده مطابقت دارند.

یک ارزیابی بالینیِ بالغ‌تر افزون بر این موارد را نیز بررسی می‌کرد:
\begin{faenum}
\faitem کالیبراسیون در کل و درون زیرگروه‌های مهم؛
\faitem بازه‌های اطمینان برای بازیابی، دقت، \lr{ROC--AUC} و \lr{PR--AUC}؛
\faitem تحلیل منحنی تصمیم یا سایر معیارهای سود خالص؛
\faitem اعتبارسنجی زمانی روی مراجعه‌های جدیدتر؛
\faitem اعتبارسنجی بیرونی در یک نظام سلامت جداگانه؛
\faitem تعریف زیرگروه‌های بالینیِ معنادار با شمار کافی رویداد مثبت.
\end{faenum}

% ============================================================
\section{محدودیت‌ها}
% ============================================================

\begin{faenum}
\faitem \textbf{نابرابری زیرگروه‌های سنی (بحرانی):} هر چهار مدل دست‌کم در یک گروه سنی همه‌ی بستری‌های مجدد زودهنگام را از دست می‌دهند \po{}\lr{FNR = 1.00}؛ شکاف ۰/۸۸۵ تا ۰/۹۶۷). گروه‌های درگیر جوان‌ترین گروه‌ها هستند و ۱ و ۶ مورد مثبت دارند، بنابراین بزرگی آن نامطمئن است، اما مدل نباید بدون اعتبارسنجی جداگانه برای این بیماران به کار رود.
\faitem \textbf{داده‌های تاریخی:} داده‌های منبع در سال ۲۰۰۸ پایان می‌یابد. مسیرهای بالینی، کدگذاری، داروها و مصرف خدمات درمانی از آن زمان به‌طور چشمگیری تغییر کرده‌اند.
\faitem \textbf{یک تقسیم نگه‌داشته‌شده:} نتایج اصلی آزمون از یک تقسیم گروهیِ بیمار-محور به دست آمده است. تقسیم مکرر و بازه‌های اطمینان انجام نشد.
\faitem \textbf{زیرگروه‌های کوچک:} چند گروه سنی و نژادی رویدادهای مثبت اندکی دارند که قابلیت اتکای برآوردهای \lr{FNR} زیرگروه‌ها را محدود می‌کند.
\faitem \textbf{ثبت پیامد:} بستری‌های مجددی که خارج از نظام بیمارستانیِ منبع رخ داده‌اند ممکن است مشاهده نشده باشند.
\faitem \textbf{زمینه‌ی بالینی گمشده:} عوامل اجتماعی تعیین‌کننده‌ی سلامت، روند تفصیلی آزمایشگاهی، شدت بیماری و ارزیابی پزشک به‌طور ناقص نمایندگی شده‌اند.
\faitem \textbf{بازنمایی ویژگی‌ها:} شناسه‌های اداریِ اسمی تا حدی عددی در نظر گرفته شدند و کدهای تشخیصی خامِ پرکاردینالیتی یک‌داغ کدگذاری شدند. این انتخاب‌ها ممکن است نوفه بیفزایند.
\faitem \textbf{داده‌ی گمشده:} برخی فیلدها داده‌ی گمشده‌ی بسیار زیاد دارند. در نظر گرفتن «گمشده» به‌عنوان یک دسته ممکن است اطلاعات را حفظ کند اما می‌تواند فرایندهای گردآوری داده را نیز کدگذاری کند.
\faitem \textbf{نبود تحلیل علّی رسمی:} ارتباط‌ها و اهمیت ویژگی‌ها روابط علّی را ثابت نمی‌کنند.
\faitem \textbf{توان تفکیک متوسط:} \lr{ROC--AUC} نزدیک ۰/۶۸ نشان‌دهنده‌ی توان تفکیک محدود است، نه ابزار پیش‌بینی بالینیِ بسیار دقیق.
\faitem \textbf{بار هشدار بالا:} آستانه‌های بازیابی‌محور مثبت‌های کاذب را به‌طور چشمگیری افزایش می‌دهند و بنابراین منابع عملیاتی می‌طلبند.
\faitem \textbf{ارزیابی ناقص زیرگروه‌ها برای مدل نهایی:} معیارهای انصاف برای مدل‌های مرحله‌ی الف محاسبه شد؛ مدل نهایی مرحله‌ی ب برای همه‌ی معیارهای زیرگروه دوباره ارزیابی نشد.
\faitem \textbf{نبود ارزیابی بالینیِ آینده‌نگر:} تحلیل نمی‌سنجد که آیا استفاده از مدل پیامدها را بهبود می‌دهد، بستری مجدد را کاهش می‌دهد یا مصرف خدمات درمانی را کم می‌کند.
\faitem \textbf{قرارداد امتیازدهی زیرگروه‌ها:} گروه‌های بدون مورد مثبت بازیابی صفر می‌گیرند \code{(zero\_division=0)} که شکاف‌های ظاهری را بزرگ می‌کند؛ \lr{FNR} برای آن گروه‌ها تعریف‌نشده است.
\faitem \textbf{دامنه‌ی انصاف:} انصاف فقط برای مدل‌های مرحله‌ی الف، فقط با بازیابی/\lr{FNR} و فقط برای نژاد، جنسیت و سن سنجیده شد؛ کالیبراسیون و برابری پیش‌بینی‌کنندگی بررسی نشدند.
\faitem \textbf{آستانه برای مدل تکمیلی:} \lr{XGBoost} مرحله‌ی ب در آستانه‌ی ثابت ۰/۵۰ ارزیابی شد؛ آستانه‌ی آن دوباره تنظیم نشد و \lr{PR--AUC} و معیارهای زیرگروه برای آن دوباره محاسبه نگردید. آستانه‌ی ۰/۴۱ به مدل مرحله‌ی الف مربوط است.
\faitem \textbf{عدم قطعیت آستانه:} آستانه روی یک تقسیم اعتبارسنجی انتخاب شد؛ پایداری آن در تقسیم‌های مختلف و بازه‌ی اطمینان بازیابیِ حاصل محاسبه نشد.
\faitem \textbf{بازنمایی:} نمونه عمدتاً قفقازی (۷۴/۷۸\%) و آفریقایی‌آمریکایی (۱۸/۸۸\%) است؛ گروه‌های دیگر کوچک‌اند.
\end{faenum}

% ============================================================
\section{پیامدهای عملی}
% ============================================================

نتایج چند پیامد محتاطانه برای یک جریان کاریِ تحلیلیِ بالینی در آینده را پشتیبانی می‌کند.

نخست، مصرف قبلی خدمات بستری را می‌توان در یک صف بازبینی پس از ترخیص وارد کرد، زیرا سیگنالی بازتولیدپذیر در تحلیل‌های توصیفی و یادگیری ماشین است. دوم، مقصد ترخیص می‌تواند در طبقه‌بندی ریسک نقش داشته باشد، اما باید همراه با زمینه‌ی بالینی تفسیر شود، نه به‌عنوان متغیر علّیِ قابل‌تغییر. سوم، آستانه‌ی کاری باید بر اساس ظرفیت موجود برای بازبینی انسانی انتخاب شود، نه با بیشینه‌کردن یک معیار عمومی.

یک جریان کاری عملی را می‌توان چنین ساخت (جهت جریان از راست به چپ است):
\begin{center}\faindent=0pt\relax
مراجعه‌ی بیمارستانی $\leftarrow$ امتیاز ریسک $\leftarrow$ آستانه‌ی مبتنی بر ظرفیت $\leftarrow$ بازبینی بالینی انسانی $\leftarrow$ پیگیری مناسب.
\end{center}
مدل نباید به‌طور خودکار درمانی را آغاز کند یا خدمتی را رد کند.

\subsection{استفاده برای اولویت‌بندی، نه خودکارسازی}

با دقتی در حدود ۱۵ تا ۱۸\%، بیشتر هشدارها کاذب‌اند؛ بنابراین مدل بیشتر برای اولویت‌بندی یک صف بازبینی دستی مناسب است تا برای راه‌اندازی مداخله‌های خودکار. آستانه یک اهرم عملیاتی است: در ۰/۴۱ مدل حدود ۵۹\% ترخیص‌ها را علامت می‌زند (۱۱,۸۴۶ از ۱۹,۹۴۸) و ۷۹/۷\% بستری‌های مجدد زودهنگام را می‌گیرد، در حالی که در آستانه‌ی پیش‌فرض ۰/۵۰ تقریباً ۳۶\% را علامت می‌زند (استخراج‌شده) و ۵۹/۸\% را می‌گیرد؛ آستانه باید بر اساس ظرفیت بازبینی و هزینه‌ی نسبیِ از دست رفتن یک بستری مجدد انتخاب شود. بیمارانی که چند بستری قبلی دارند به‌طور مداوم برجسته‌اند، بنابراین منابع پیگیری (تماس‌های هماهنگی مراقبت، ارجاع به مراقبت در منزل، نوبت‌های پیگیریِ زودهنگام) می‌تواند برای این گروه در اولویت قرار گیرد؛ داده‌ها ارتباط را نشان می‌دهند، نه اینکه این اقدام‌ها بستری مجدد را کاهش دهند.

\subsection{شرایطی که در آن نباید از مدل استفاده کرد}

مدل نباید بدون پرداختن به مسئله‌ی سن به کار رود: هر چهار مدل دست‌کم در یک گروه سنی همه‌ی بستری‌های مجدد زودهنگام را از دست می‌دهند. بیماران زیر ۲۰ سال باید کنار گذاشته یا جداگانه اعتبارسنجی شوند، بازیابی در گروه‌های سنی باید با بازه‌های اطمینان گزارش شود، و پیش از هر پایلوت باید کالیبراسیون یا آستانه‌های آگاه از سن در نظر گرفته شود. سرانجام، چون داده‌ها در سال ۲۰۰۸ پایان می‌یابند، هر استقرار واقعی به اعتبارسنجی دوباره روی داده‌ی امروزی نیاز دارد.

پیش از پیاده‌سازی، جریان کار باید به‌صورت آینده‌نگر آزمایش شود. یک پایلوت مفید حجم هشدار، پذیرش پزشکان، زمان بازبینی، تکمیل پیگیری، پیامدهای بستری مجدد، کالیبراسیون و عملکرد زیرگروه‌ها را اندازه می‌گرفت. هر تغییری در پیامد بالینی باید با طرح مطالعه‌ی مناسب ارزیابی شود، نه اینکه از عملکرد پیش‌بینی استنباط گردد.

% ============================================================
\section{بازتولیدپذیری و محیط محاسباتی}
% ============================================================

پروژه با پایتون پیاده‌سازی شد: \lr{pandas} و \lr{NumPy} برای پردازش داده؛ \lr{SciPy} برای محاسبات آماری؛ \lr{scikit-learn} برای پیش‌پردازش، تقسیم‌بندی، مدل‌سازی، اعتبارسنجی متقاطع و معیارها؛ \lr{XGBoost} برای تقویت گرادیانی؛ و \lr{Matplotlib}/\lr{Seaborn} برای مصورسازی.

\begin{table}[H]
\centering
\caption{محیط محاسباتی و اجزای تحلیلی.}
\label{tab:software}
{\faindent=0pt\relax\begin{tabularx}{\textwidth}{>{\raggedleft\arraybackslash}X>{\beginR}r<{\endR}}
\toprule
ابزارها & دسته\\

\midrule
\lr{Python 3.13.5} در محیط اصلی گزارش شده است & برنامه‌نویسی\\
\lr{pandas, NumPy} & پردازش داده\\
\lr{SciPy}، روش‌های جدول پیشامد، اندازه‌های اثر & آمار\\
\lr{scikit-learn, XGBoost} & یادگیری ماشین\\
\code{GroupShuffleSplit}، \code{StratifiedGroupKFold}، \code{RandomizedSearchCV} & اعتبارسنجی\\
\lr{Matplotlib, Seaborn} & مصورسازی\\
\lr{Jupyter Notebook / Anaconda} & محیط\\

\bottomrule
\end{tabularx}}
\end{table}

پروژه‌ی اصلی از یک مسیر نسبی برای داده استفاده می‌کند و شکل‌ها را در پوشه‌ی \code{Figures/} ذخیره می‌کند. بازتولید عددی دقیق ممکن است با نسخه‌ی کتابخانه‌ها اندکی تفاوت کند، اما بذر اصلی تقسیم ۱۲۳ است.

بذر تصادفی برای تقسیم آموزش/آزمون، تقسیم اعتبارسنجی، مدل‌ها و جست‌وجوی ابرپارامتر ۱۲۳ است. نسخه‌ی بسته‌ها در پروژه‌ی اصلی ثبت نشد و هنگام ثابت‌کردن محیط باید قفل شود. دفترچه داده را با مسیر نسبیِ سبک ویندوز می‌خواند که در مک‌اواس یا لینوکس باید تغییر کند، و باید از درون پوشه‌ی خودش اجرا شود تا مسیرهای نسبی درست حل شوند. نقاط بازرسیِ مورد انتظار عبارت‌اند از جدول داده‌ی خامِ $101{,}766\times 50$، هم‌گروه مدل‌سازی با ۹۹,۳۴۳ مراجعه، تقسیم آموزش/آزمون با ۷۹,۳۹۵ و ۱۹,۹۴۸ مراجعه، و هم‌پوشانی صفر بیمار.

یک داشبورد \lr{HTML} تعاملی \po{}\lr{Chart.js}، با پالایه‌های سن و جنسیت) همراه پروژه است. این داشبورد توزیع پیامد دودویی، دو نمودار آستانه‌ی \lr{XGBoost} (۰/۴۱ و پیش‌فرض ۰/۵۰، ارزیابی‌شده روی مجموعه‌ی آزمون نگه‌داشته‌شده و بنابراین بی‌اثر از پالایه‌ها) و یک کارت هشدار انصاف سنی را نشان می‌دهد که شکاف‌های \lr{FNR} نژاد، جنسیت و سن را برای هر چهار مدل خلاصه می‌کند.

% ============================================================
\section{نتیجه‌گیری}
% ============================================================

این پژوهش یک گردش کار کامل علم داده‌های پزشکی را برای پیش‌بینی بستری مجدد ظرف ۳۰ روز در بیماران بستری مبتلا به دیابت ارائه می‌دهد. تحلیل، بازبینی کیفیت داده، اپیدمیولوژی توصیفی، اندازه‌های اثر تفسیرپذیر، یادگیری ماشین آگاه از نشت، بهینه‌سازی آستانه و تحلیل خطای زیرگروه‌ها را با هم ترکیب می‌کند.

قوی‌ترین سیگنال تکرارشونده مصرف قبلی خدمات بستری است، در حالی که مقصد ترخیص بزرگ‌ترین ارتباط طبقه‌ای و یک ویژگی برجسته‌ی مدل را فراهم می‌کند. با این حال، بیشتر ارتباط‌های تکی کوچک‌اند و مدل‌های پیش‌بینی تنها توان تفکیک متوسطی دارند. \lr{XGBoost} به \lr{ROC--AUC} حدود ۰/۶۸ رسید، و آستانه‌گذاریِ بازیابی‌محور حساسیت را به‌طور چشمگیری افزایش داد، به بهای بار بزرگ مثبت‌های کاذب.

این مطالعه همچنین نشان می‌دهد چرا پیش‌بینی در سلامت نباید به یک عدد صحت فروکاسته شود. تقسیم‌بندی در سطح بیمار از یک شکل مهم نشت جلوگیری می‌کند، \lr{PR--AUC} در عدم توازن کلاس‌ها آگاهی‌بخش است، و معیارهای زیرگروه به زمینه‌ی اندازه‌ی نمونه نیاز دارند. نتایج انصاف گروه‌های سنی نشان می‌دهد وقتی موارد مثبت نادرند، برآوردهای ظاهراً چشمگیرِ نابرابری چقدر می‌توانند ناپایدار باشند.

بنابراین تفسیر مناسب محافظه‌کارانه است: مدل می‌تواند از اولویت‌بندی و بازبینی انسانی پشتیبانی کند، اما نباید تصمیم‌گیرنده‌ی بالینیِ خودمختار تلقی شود. اعتبارسنجی بیرونیِ معاصر، کالیبراسیون، کمّی‌سازی عدم قطعیت، ارزیابی زیرگروه‌ها و سنجش آینده‌نگر اثر بالینی، گام‌های لازم پیش از هر کاربرد واقعی است.

\section*{دسترس‌پذیری داده}
تحلیل از مجموعه‌داده‌ی عمومی «دیابت، ۱۳۰ بیمارستان آمریکا، سال‌های ۱۹۹۹ تا ۲۰۰۸» استفاده می‌کند که از طریق مخزن یادگیری ماشین \lr{UCI} با مجوز \lr{CC BY 4.0} منتشر شده است. آرشیو پروژه‌ی ارائه‌شده شامل فایل‌های داده، دفترچه و شکل‌های تولیدشده‌ی مورد استفاده برای این دست‌نوشته است.

\section*{بیانیه‌ی اخلاق و کاربرد بالینی}
این یک تحلیل ثانویه روی مجموعه‌داده‌ی عمومیِ بدون شناسه است. این دست‌نوشته هیچ توصیه‌ی بالینی در سطح فرد ارائه نمی‌کند. نتایج پیش‌بینی برای تحلیل روش‌شناختی و آموزشی ارائه شده‌اند و بدون اعتبارسنجی بالینیِ مستقل و حکمرانی مناسب نباید مستقیماً در مراقبت از بیمار به کار روند.

% ============================================================
\appendix
\section{تعداد دقیق خطاهای مدل‌ها}
% ============================================================

\begin{table}[H]
\centering
\caption{تعداد خطاها در آستانه‌های بازیابی‌محور مرحله‌ی الف.}
\label{tab:errorcounts}
{\faindent=0pt\relax\begin{tabular}{>{\beginR}r<{\endR}>{\beginR}r<{\endR}>{\beginR}r<{\endR}>{\beginR}r<{\endR}>{\beginR}r<{\endR}>{\beginR}r<{\endR}}
\toprule
\lr{FNR} & \lr{TN} & \lr{FP} & \lr{FN} & \lr{TP} & مدل\\

\midrule
۰/۲۰۹ & ۷,۱۴۹ & ۱۰,۵۳۹ & ۴۷۳ & ۱,۷۸۷ & رگرسیون لجستیک\\
۰/۲۱۴ & ۷,۰۸۰ & ۱۰,۶۰۸ & ۴۸۴ & ۱,۷۷۶ & درخت تصمیم\\
۰/۱۷۰ & ۶,۶۱۹ & ۱۱,۰۶۹ & ۳۸۴ & ۱,۸۷۶ & جنگل تصادفی\\
۰/۲۰۳ & ۷,۶۴۴ & ۱۰,۰۴۴ & ۴۵۸ & ۱,۸۰۲ & \lr{XGBoost}\\

\bottomrule
\end{tabular}}
\end{table}

\section{ممیزی سازگاری و تفسیر}

پروژه‌ی ارائه‌شده خروجی‌های متعددی دارد که مراحل تحلیلیِ متفاوتی را نمایندگی می‌کنند. دست‌نوشته‌ی حاضر عمداً آن‌ها را از هم تفکیک می‌کند، نه اینکه یک عدد واحد را به همه‌ی بخش‌ها تحمیل کند.

\begin{table}[H]
\centering
\caption{تمایزهای تحلیلیِ کلیدی که در این دست‌نوشته حفظ شده‌اند.}
\label{tab:audit}
{\faindent=0pt\relax\begin{tabularx}{\textwidth}{>{\raggedleft\arraybackslash}X>{\raggedleft\arraybackslash}X>{\beginR}r<{\endR}}
\toprule
مرحله‌ی ب / نهایی & مرحله‌ی الف & موضوع\\

\midrule
پیکربندی انتخاب‌شده با \code{RandomizedSearchCV} & پیکربندی ثابت ۲۵۰ درختی & پارامترهای \lr{XGBoost}\\
ثابت روی ۰/۵۰ & تنظیم‌شده روی اعتبارسنجی؛ ۰/۴۱ برای \lr{XGBoost} & آستانه\\
۰/۶۱۹ & ۰/۷۹۷ & بازیابی \lr{XGBoost}\\
۰/۱۸۱ & ۰/۱۵۲ & دقت \lr{XGBoost}\\
تقریباً ۰/۶۸ & ۰/۶۸۳ & \lr{ROC--AUC} مدل \lr{XGBoost}\\
گزارش مدل نهایی جست‌وجوشده & مقایسه‌ی نقاط کاری & هدف\\
به‌طور کامل دوباره محاسبه نشده & محاسبه شده & معیارهای انصاف\\

\bottomrule
\end{tabularx}}
\end{table}

این تمایز مانع می‌شود خواننده دو پیکربندی متفاوت مدل را یک آزمایش واحد بپندارد.

\subsection*{ناسازگاری‌های میان شرح دفترچه و خروجی‌های اجراشده}

در طول آماده‌سازی، توضیحات نوشتاریِ دفترچه با خروجی‌های اجراشده‌اش مقایسه شد. خروجی‌های اجراشده (نتایج ذخیره‌شده‌ی سلول‌ها) قابل‌اعتمادترین منبع تلقی شدند.

{\small
\begin{longtable}{>{\raggedleft\arraybackslash}p{6.4cm}>{\raggedleft\arraybackslash}p{8.3cm}>{\beginR}r<{\endR}}
\caption{ممیزی سازگاریِ دفترچه‌ی منبع.}\label{tab:audit2}\\
\toprule
راه‌حل به‌کاررفته & مسئله و آنچه خروجی اجراشده نشان می‌دهد & ردیف\\
\midrule
\endfirsthead
\toprule
راه‌حل به‌کاررفته & مسئله و آنچه خروجی اجراشده نشان می‌دهد & ردیف\\
\midrule
\endhead
\lr{XGBoost} به دلیل بهترین \lr{ROC--AUC}، \lr{PR--AUC}، صحت متوازن و \lr{F1} توجیه می‌شود، نه به دلیل بیشترین بازیابی. & متن «مدل منتخب - \lr{XGBoost}» می‌گوید \lr{XGBoost} در آستانه‌ی ۰/۵۰ بیشترین بازیابی (۰/۸۰۰) و \lr{ROC--AUC} برابر ۰/۶۸۴ را داشت. در ۰/۵۰، بازیابی \lr{XGBoost} برابر ۰/۵۹۸ است و درخت تصمیم بیشترین (۰/۶۵۹) را دارد؛ پس از تنظیم، جنگل تصادفی بیشترین بازیابی (۰/۸۳۰) را دارد؛ \lr{ROC--AUC} مدل \lr{XGBoost} برابر ۰/۶۸۳ است. هیچ خروجی اجراشده‌ای بازیابی ۰/۸۰۰ را برای \lr{XGBoost} نشان نمی‌دهد (ممکن است از اجرای قبلی باشد). & ۱\\
به‌عنوان مرحله‌ی الف (ابرپارامترهای ثابت، آستانه‌ی تنظیم‌شده‌ی ۰/۴۱) و مرحله‌ی ب (ابرپارامترهای جست‌وجوشده، آستانه‌ی ۰/۵۰، تکمیلی) گزارش شد. & شرح پرسش ۹ مدل «قهرمان» \lr{XGBoost} را در آستانه‌ی ۰/۴۱ با بازیابی ۰/۷۹۷ توصیف می‌کند، اما سلول‌های پایانی \lr{XGBoost} را با \code{RandomizedSearchCV} دوباره برازش می‌دهند و در آستانه‌ی ۰/۵۰ ارزیابی می‌کنند (بازیابی ۰/۶۱۹). & ۲\\
نابرابری سنی به‌عنوان یافته‌ی بحرانی همراه با عدم قطعیت نمونه‌ی کوچک گزارش شد؛ صفرِ جنسیت به‌عنوان مصنوع داده توضیح داده شد. & شرح انصاف از «سوگیری الگوریتمیِ شدید» بر حسب سن و \code{min\_Recall = 0.000} برای جنسیت سخن می‌گوید. شکاف \lr{FNR} سن ۰/۸۸۵ تا ۰/۹۶۷ است و در هر چهار مدل دست‌کم در یک گروه \lr{FNR = 1.00} است؛ آن گروه‌ها ۱ و ۶ مورد مثبت دارند. صفر جنسیت از یک گروه «گمشده»ی تک‌مراجعه‌ای بدون مورد مثبت ناشی می‌شود. انصاف فقط برای مدل‌های مرحله‌ی الف محاسبه شد. & ۳\\
هر دو گزارش و بر حسب مرحله برچسب‌گذاری شدند. & دو اندازه‌ی هم‌گروه و شیوع ظاهر می‌شود: تحلیل اکتشافی ۱۰۱,۷۶۶ مراجعه (۱۱/۱۶\%) دارد؛ مدل‌سازی پس از حذف ۲,۴۲۳ ترخیص فوت/هاسپیس ۹۹,۳۴۳ مراجعه (۱۱/۳۹\%) دارد. & ۴\\
۰/۶۸۳ به‌کار رفت (تفاوت گردکردن). & متن دفترچه می‌گوید \lr{PR--AUC} مدل \lr{XGBoost} برابر ۰/۲۳۸ و \lr{ROC--AUC} برابر ۰/۶۸۴ است؛ خروجی ۰/۲۳۸ و ۰/۶۸۳ را نشان می‌دهد. & ۵\\
فقط \code{project-04.ipynb} به‌عنوان منبع مرجع به کار رفت. & عنوان و نام فایل دفترچه‌ی اصلی با نسخه‌های نقطه‌ی بازرسی‌اش تفاوت دارد؛ پوشه‌ی نقطه‌ی بازرسی یک نسخه‌ی قدیمی‌تر با نام متفاوت دارد. & ۶\\
\bottomrule
\end{longtable}
}

\section{واژه‌نامه‌ی اصطلاحات آماری}

برای اطمینان از یکدستی ترجمه، معادل فارسی اصطلاحات آماری و یادگیری ماشین که در این مقاله به کار رفته‌اند در جدول زیر گردآوری شده است.

\begin{table}[H]
\centering
\caption{واژه‌نامه‌ی اصطلاحات آماری (انگلیسی به فارسی).}
\label{tab:glossary}
\small
{\faindent=0pt\relax\begin{tabularx}{\textwidth}{>{\beginR}r<{\endR}>{\raggedleft\arraybackslash}X}
\toprule
معادل فارسی & اصطلاح انگلیسی\\
\midrule
صحت & \lr{Accuracy}\\
دقت (ارزش اخباری مثبت) & \lr{Precision (PPV)}\\
بازیابی (حساسیت) & \lr{Recall / Sensitivity}\\
نرخ منفی کاذب & \lr{False Negative Rate (FNR)}\\
صحت متوازن & \lr{Balanced Accuracy}\\
ماتریس درهم‌ریختگی & \lr{Confusion Matrix}\\
منحنی مشخصه‌ی عملکرد گیرنده / مساحت زیر منحنی & \lr{ROC Curve / AUC}\\
منحنی دقت--بازیابی & \lr{Precision--Recall Curve}\\
عدم توازن کلاس‌ها & \lr{Class Imbalance}\\
تفاوت میانگین استانداردشده & \lr{Standardized Mean Difference (SMD)}\\
اندازه‌ی اثر & \lr{Effect Size}\\
ضریب \lr{V} کرامر & \lr{Cram\'er's V}\\
همبستگی پیرسون & \lr{Pearson Correlation}\\
دامنه‌ی بین‌چارکی & \lr{Interquartile Range (IQR)}\\
داده‌ی پرت & \lr{Outlier}\\
جای‌گذاری مقادیر گمشده & \lr{Imputation}\\
اعتبارسنجی متقاطع & \lr{Cross-Validation}\\
ابرپارامتر & \lr{Hyperparameter}\\
آستانه‌ی تصمیم & \lr{Decision Threshold}\\
کالیبراسیون (واسنجی) & \lr{Calibration}\\
بازه‌ی اطمینان & \lr{Confidence Interval}\\
نشت داده & \lr{Data Leakage}\\
هم‌گروه & \lr{Cohort}\\
مجموعه‌ی آزمون نگه‌داشته‌شده & \lr{Held-out Test Set}\\
تحلیل زیرگروه‌ها & \lr{Subgroup Analysis}\\
کدگذاری یک‌داغ & \lr{One-hot Encoding}\\
صفر-متورم & \lr{Zero-inflated}\\
نسبت افزایش & \lr{Lift}\\
رگرسیون لجستیک & \lr{Logistic Regression}\\
درخت تصمیم & \lr{Decision Tree}\\
جنگل تصادفی & \lr{Random Forest}\\
تقویت گرادیانی & \lr{Gradient Boosting}\\
\bottomrule
\end{tabularx}}
\end{table}

% ============================================================
% مراجع
% ============================================================
\clearpage
\section*{مراجع}
{\small\begingroup\faindent=0pt\relax
\newcommand{\refitem}[2]{\par\vskip.4em\noindent\hbox to 1.4cm{\hfil[#1]}\hskip.6em\begin{minipage}[t]{\dimexpr\textwidth-2.2cm\relax}\raggedright\latinfont\beginL #2\endL\end{minipage}}
\refitem{1}{B. Strack, J. P. DeShazo, C. Gennings, J. L. Olmo, S. Ventura, K. J. Cios, and J. N. Clore, ``Impact of HbA1c measurement on hospital readmission rates: Analysis of 70,000 clinical database patient records,'' \textit{BioMed Research International}, vol. 2014, Article 781670, 2014. doi:10.1155/2014/781670.}
\refitem{2}{J. Clore, K. Cios, J. DeShazo, and B. Strack, \textit{Diabetes 130-US Hospitals for Years 1999--2008}. UCI Machine Learning Repository, 2014. Available at: \url{https://archive.ics.uci.edu/dataset/296/diabetes-130-us-hospitals-for-years-1999-2008}.}
\refitem{3}{T. Chen and C. Guestrin, ``XGBoost: A scalable tree boosting system,'' in \textit{Proceedings of the 22nd ACM SIGKDD International Conference on Knowledge Discovery and Data Mining}, pp. 785--794, 2016. doi:10.1145/2939672.2939785.}
\refitem{4}{F. Pedregosa, G. Varoquaux, A. Gramfort, et al., ``Scikit-learn: Machine learning in Python,'' \textit{Journal of Machine Learning Research}, vol. 12, pp. 2825--2830, 2011.}
\refitem{5}{C. R. Harris, K. J. Millman, S. J. van der Walt, et al., ``Array programming with NumPy,'' \textit{Nature}, vol. 585, pp. 357--362, 2020. doi:10.1038/s41586-020-2649-2.}
\refitem{6}{P. Virtanen, R. Gommers, T. E. Oliphant, et al., ``SciPy 1.0: Fundamental algorithms for scientific computing in Python,'' \textit{Nature Methods}, vol. 17, pp. 261--272, 2020. doi:10.1038/s41592-019-0686-2.}
\refitem{7}{J. D. Hunter, ``Matplotlib: A 2D graphics environment,'' \textit{Computing in Science \& Engineering}, vol. 9, no. 3, pp. 90--95, 2007. doi:10.1109/MCSE.2007.55.}
\refitem{8}{M. L. Waskom, ``Seaborn: Statistical data visualization,'' \textit{Journal of Open Source Software}, vol. 6, no. 60, p. 3021, 2021. doi:10.21105/joss.03021.}
\refitem{9}{J. Cohen, \textit{Statistical Power Analysis for the Behavioral Sciences}, 2nd ed. Lawrence Erlbaum Associates, 1988.}
\refitem{10}{H. Cram\'er, \textit{Mathematical Methods of Statistics}. Princeton University Press, 1946.}
\endgroup}

\end{document}
